\chapter{Polynomials and rational fractions}

Throughout this chapter A denotes a commutative ring.

\section{POLYNOMIALS}

\subsection{Definition of polynomials}

Let I be a set. We recall (III, p.~452) that the free commutative algebra on I over A is denoted by \(\mathbf{A}[(\mathbf{X}_i)_{i\in I}]\) or \(\mathbf{A}[\mathbf{X}_i]_{i\in I}\) . The elements of this algebra are called polynomials with respect to the indeterminates \(\mathbf{X}_i\) (or in the indeterminates \(\mathbf{X}_i\) ) with coefficients in A. Let us recall that the indeterminate \(\mathbf{X}_i\) is the canonical image of \(i\) in the free commutative algebra on I over A; sometimes it is convenient to denote this image by another symbol such as \(\mathbf{X}_i'\) , \(\mathbf{Y}_i\) , \(\mathbf{T}_i\) , etc. This convention is often introduced by a phrase such as: « Let \(\mathbf{Y} = (\mathbf{Y}_i)_{i\in I}\) be a family of indeterminates »; in this case the algebra of polynomials in question is denoted by \(\mathbf{A}[\mathbf{Y}]\) . When \(\mathrm{I} = \{1,2,\dots,n\}\) , one writes \(\mathbf{A}[\mathbf{X}_1,\mathbf{X}_2,\dots,\mathbf{X}_n]\) in place of \(\mathbf{A}[(\mathbf{X}_i)_{i\in I}]\) .

For \(\nu \in \mathbf{N}^{(I)}\) we put

\[
\mathbf {X} ^ {\nu} = \prod_ {i \in I} \mathbf {X} _ {i} ^ {\nu_ {i}}.
\]

Then \((\mathbf{X}^{\nu})_{\nu \in \mathbf{N}^{(I)}}\) is a basis of the A-module \(\mathbf{A}[(\mathbf{X}_i)_{i\in I}]\) . The \(\mathbf{X}^{\nu}\) are called monomials in the indeterminates \(\mathbf{X}_i\) . For \(\nu = 0\) we obtain the unit element of \(\mathbf{A}[(\mathbf{X}_i)_{i\in I}]\) . Every polynomial \(u\in \mathbf{A}[(\mathbf{X}_i)_{i\in I}]\) can be written in exactly one way in the form

\[
u = \sum_ {\nu \in \mathbf {N} ^ {(I)}} \alpha_ {\nu} X ^ {\nu}
\]

where \(\alpha_{\nu}\in A\) and the \(\alpha_{\nu}\) are zero except for a finite number; the \(\alpha_{\nu}\) are called the coefficients of u; the \(\alpha_{\nu}X^{\nu}\) are called the terms of u (often the element \(\alpha_{\nu}X^{\nu}\) is called the term in \(X^{\nu}\) ), in particular the term of \(\alpha_{0}X^{0}\) , identified with \(\alpha_{0}\) , is called the constant term of u. When \(\alpha_{\nu}=0\) , we say by abuse of language that u contains no element in \(X^{\nu}\) ; in particular when \(\alpha_{0}=0\) , we say that u is a polynomial without constant term (III, p.~453). Any scalar multiple of 1 is called a constant polynomial.

Let B be a commutative ring and \(\rho: \mathbf{A} \to \mathbf{B}\) a ring homomorphism. We consider \(\mathrm{B}[(\mathbf{X}_i)_{i \in I}]\) as an A-algebra by means of \(\rho\) . Thus the mapping \(\sigma\) of \(\mathbf{A}[(\mathbf{X}_i)_{i \in I}]\) into \(\mathrm{B}[(\mathbf{X}_i)_{i \in I}]\) which transforms \(\sum \alpha_\nu \mathbf{X}^\nu\) into \(\sum \rho(\alpha_\nu) \mathbf{X}^\nu\) is a homomorphism of A-algebras; if \(u \in \mathbf{A}[(\mathbf{X}_i)_{i \in I}]\) , we sometimes denote by \(^p u\) the image of \(u\) by this homomorphism. The homomorphism of \(\mathbf{B} \otimes_{\mathbf{A}} \mathbf{A}[(\mathbf{X}_i)_{i \in I}]\) into \(\mathrm{B}[(\mathbf{X}_i)_{i \in I}]\) canonically defined by \(\sigma\) transforms, for every \(i \in I\) , \(1 \otimes \mathbf{X}_i\) into \(\mathbf{X}_i\) ; this is an isomorphism of B-algebras (III, p.~449).

Let M be a free A-module with basis \((e_{i})_{i\in I}\) . There exists precisely one unital homomorphism \(\varphi\) of the symmetric algebra \(\mathbf{S}(\mathbf{M})\) into the algebra \(\mathbf{A}[(\mathbf{X}_{i})_{i\in I}]\) such that \(\varphi(e_{i})=\mathbf{X}_{i}\) for each \(i\in I\) , and this homomorphism is an isomorphism (III, p.~506). This isomorphism is said to be canonical. It allows us to apply to polynomial algebras certain properties of symmetric algebras. For example, let \((\mathrm{I}_{\lambda})_{\lambda\in L}\) be a partition of I. Let \(\varphi_{\lambda}\) be the homomorphism of \(P_{\lambda}=\mathbf{A}[(\mathbf{X}_{i})_{i\in I_{\lambda}}]\) into \(P=\mathbf{A}[(\mathbf{X}_{i})_{i\in I}]\) which transforms \(X_{i}\) (qua element of \(P_{\lambda}\) ) into \(X_{i}\) (qua element of P). Then the \(\varphi_{\lambda}\) define a homomorphism of the algebra \(\otimes_{\lambda\in L}P_{\lambda}\) into the algebra P,

and this homomorphism is an isomorphism (III, p.~503, Prop. 9).

Let \(\mathbf{E}\) be an \(\mathbf{A}\) -module, and put \(\mathbf{E} \otimes_{\mathbf{A}} \mathbf{A}[(\mathbf{X}_i)_{i \in I}] = \mathbf{E}[(\mathbf{X}_i)_{i \in I}]\) . The elements of the \(\mathbf{A}\) -module \(\mathbf{E}[(\mathbf{X}_i)_{i \in I}]\) are called polynomials in the indeterminates \(\mathbf{X}_i\) with coefficients in \(\mathbf{E}\) . Such a polynomial can be written in just one way as \(\sum_{\nu \in \mathbf{N}^{(I)}} e_\nu \otimes \mathbf{X}^\nu\) , where \(e_\nu \in \mathbf{E}\) and the \(e_\nu\) are zero for all but a finite number of

suffixes; we frequently write \(e_{\nu}X^{\nu}\) instead of \(e_{\nu}\otimes X^{\nu}\) .

\subsection{Degrees}

Let \(\mathbf{P} = \mathbf{A}[(\mathbf{X}_i)_{i\in I}]\) be a polynomial algebra. For each integer \(n\in \mathbf{N}\) let \(\mathbf{P}_n\) be the submodule of \(\mathbf{P}\) generated by the monomials \(\mathbf{X}^{\nu}\) such that \(|\nu| = \sum_{i\in I}\nu_i\)

equals \(n\) . Then \((\mathbf{P}_n)_{n \in \mathbf{N}}\) is a graduation which turns \(\mathbf{A}[(\mathbf{X}_i)_{i \in I}]\) into a graded algebra of type \(\mathbf{N}\) (III, p.~459). The homogeneous elements of degree \(n\) in \(\mathbf{A}[(\mathbf{X}_i)_{i \in I}]\) are sometimes called forms of degree \(n\) with respect to the indeterminates \(\mathbf{X}_i\) .

When we are dealing with the degree of inhomogeneous polynomials, we shall agree to adjoin to the set N of natural numbers, an element written \(-\infty\) and to extend the order relation and the addition of N to \(N \cup \{-\infty\}\) by the following conventions, where \(n \in N\) ,

\[
- \infty <   n, (- \infty) + n = n + (- \infty) = - \infty , (- \infty) + (- \infty) = - \infty .
\]

Let \(u = \sum_{\nu \in \mathbf{N}^{(I)}} \alpha_{\nu} X^{\nu}\) be a polynomial. The homogeneous component \(u_{n}\) of degree

\(n\) of \(u\) (for the graduation of type \(\mathbf{N}\) defined above) is equal to \(\sum_{|\nu| = n} \alpha_{\nu} X^{\nu}\) , and we

clearly have \(u = \sum_{n \in \mathbb{N}} u_n\) . If \(u \neq 0\) , the \(u_n\) are not all zero, and we define the degree

(or total degree) of u, written \(\deg u\) , as the greatest of the numbers n such that \(u_{n} \neq 0\) ; in other words (III, p.~453), the degree of u is the largest of the integers \(|\nu|\) for the multi-indices \(\nu\) such that \(\alpha_{\nu} \neq 0\) . When u = 0, the degree of u is \(-\infty\) by convention. For every integer \(p \in N\) , the relation \(\deg u \leqslant p\) is thus equivalent to \(\alpha_{\nu} = 0\) for every multi-index \(\nu\) with \(|\nu| > p\) ; the set of polynomials u such that \(\deg u \leqslant p\) is thus an A-submodule of \(\mathbf{A}[(\mathbf{X}_{i})_{i \in I}]\) , equal to \(P_{0} + P_{1} + \cdots + P_{p}\) with the above notations.

Let E be an A-module. The family \((\mathrm{E} \otimes \mathrm{P}_{n})_{n \in \mathbb{N}}\) is a graduation of type N of the module \(\mathrm{E}[(\mathrm{X}_{i})_{i \in I}] = \mathrm{E} \otimes_{\mathbb{A}} \mathrm{A}[(\mathrm{X}_{i})_{i \in I}]\) of polynomials with coefficients in E. We extend the conventions adopted above for the degree of inhomogeneous polynomials to this case.

PROPOSITION 1. --- Let u and v be two polynomials.

\begin{enumerate}
\def\labelenumi{(\roman{enumi})}
\tightlist
\item
  If \(\deg u \neq \deg v\) , we have
\end{enumerate}

\[
u + v \neq 0 \quad a n d \quad \deg (u + v) = \sup (\deg u, \deg v).
\]

If \(\deg u = \deg v\) we have \(\deg (u + v) \leqslant \deg u\) .

\begin{enumerate}
\def\labelenumi{(\roman{enumi})}
\setcounter{enumi}{1}
\tightlist
\item
  We have \(\deg(uv) \leqslant \deg u + \deg v\) .
\end{enumerate}

The proof is immediate.

Let \(J \subset I\) and \(B = A[(X_i)_{i \in I - J}]\) ; we shall identify \(A[(X_i)_{i \in I}]\) with \(B[(X_i)_{i \in J}]\) (III, p.~453 f.). The degree of \(u \in A[(X_i)_{i \in I}]\) , qua element of \(B[(X_i)_{i \in J}]\) , is called the degree of \(u\) with respect to the \(X_i\) of index \(i \in J\) (III, p.~454).

Let \(u = \sum_{k=0}^{n} \alpha_k X^k \in A[X]\) be a non-zero polynomial of degree \(n\) in a single

indeterminate. The coefficient \(\alpha_{n}\) which is \(\neq0\) by hypothesis, is called the leading coefficient of u. A polynomial \(u\neq0\) whose leading coefficient is equal to 1 is called a monic polynomial.

In \(\mathbf{A}[\mathbf{X}_1, \mathbf{X}_2, \ldots, \mathbf{X}_q]\) the number of monomials of total degree \(p\) is equal to the number of elements \((n_k)_{1 \leqslant k \leqslant q}\) of \(\mathbf{N}^q\) such that \(\sum_{k=1}^{q} n_k = p\) , that is \(\binom{q + p - 1}{p}\)

(Sets III, Prop. 15, p.~182).

More generally, let \(\Delta\) be a commutative monoid and \((\delta_i)_{i \in I}\) a family of elements of \(\Delta\) . There exists a unique graduation of type \(\Delta\) of the algebra \(\mathbf{A}[(\mathbf{X}_i)_{i \in I}]\) such that each monomial \(\mathbf{X}^\nu\) is of degree \(\sum_{i \in I} \nu_i \delta_i\) (III, p.~458,

example 3). The case considered above is that where \(\Delta = \mathbf{N}\) and \(\delta_{i} = 1\) . In the general case, to avoid confusion, we shall use the word « weight » instead of « degree » and « isobaric » instead of « homogeneous ». For example, there exists a unique graduation of type \(\mathbf{N}\) of the algebra \(\mathbf{A}[(\mathbf{X}_{i})_{i\geqslant 1}]\) such that \(\mathbf{X}_i\) is of weight \(i\) for each integer \(i \geqslant 1\) . The isobaric elements of weight \(n\) are the polynomials of the form \(\sum_{\nu} a_{\nu} X^{\nu}\) , where \(a_{\nu} = 0\) for \(\sum_{i \geqslant 1} i \cdot \nu_i \neq n\) .

\subsection{Substitutions}

Let E be a unital associative algebra over A and \(\mathbf{x} = (x_{i})_{i \in I}\) a family of pairwise permutable elements of E. Let \(\mathbf{X} = (\mathbf{X}_{i})_{i \in I}\) be a family of indeterminates. By III, Prop. 7, p.~449, there exists a unique unital homomorphism f of A{[}X{]} into E such that \(f(\mathbf{X}_{i}) = x_{i}\) for every \(i \in I\) . The image of an element u of A{[}X{]} under f is written \(u(\mathbf{x})\) and is called the element of E obtained by substituting \(x_{i}\) for \(X_{i}\) in u, or also the value of u for \(X_{i} = x_{i}\) . In particular, \(u = u((\mathbf{X}_{i})_{i \in I})\) . If \(I = \{1, \ldots, n\}\) we write \(u(x_{1}, \ldots, x_{n})\) in place of \(u((x_{i})_{i \in I})\) . More generally, if M is an A-module and if v is an element of

\[
\mathbf {M} [ (\mathbf {X} _ {i}) _ {i \in I} ] = \mathbf {M} \otimes_ {\mathbf {A}} \mathbf {A} [ (\mathbf {X} _ {i}) _ {i \in I} ],
\]

we denote the image of \(v\) in \(\mathbf{M} \otimes_{\mathbf{A}} \mathbf{E} = \mathbf{M}_{(\mathbf{E})}\) under the mapping \(1_{\mathbf{M}} \otimes f\) by \(v(\mathbf{x})\) .

If the homomorphism \(u \mapsto u(\mathbf{x})\) of \(\mathbf{A}[\mathbf{X}]\) into \(\mathbf{E}\) is injective, we say that the family \(\mathbf{x}\) is algebraically free over \(\mathbf{A}\) , or that the \(x_i\) are algebraically independent over \(\mathbf{A}\) . This also means that the monomials \(\mathbf{x}^\nu\) ( \(\nu \in \mathbf{N}^{(I)}\) ) are linearly independent over \(\mathbf{A}\) .

If \(\lambda\) is a unital homomorphism of \(E\) into a unital associative A-algebra \(E'\) , we have

\[
\lambda \left(u \left(\left(x _ {i}\right) _ {i \in I}\right)\right) = u \left(\left(\lambda \left(x _ {i}\right) _ {i \in I}\right)\right),\tag{1}
\]

because \(\lambda \circ f\) is a homomorphism of \(\mathbf{A}[\mathbf{X}]\) into \(\mathbf{E}'\) which maps \(\mathbf{X}_i\) to \(\lambda(x_i)\) .

Let \(u \in A[X]\) . If E is commutative, the mapping \(x \mapsto u(x)\) of \(E^{I}\) into E is called the polynomial function defined by \(u\) (and the algebra E); we shall sometimes denote it by \(\tilde{u}\) (or even just \(u\) ).

Let \(\mathbf{Y} = (\mathbf{Y}_j)_{j\in J}\) be another family of indeterminates, and let us take for E the polynomial algebra \(\mathbf{A}[\mathbf{Y}]\) . Given \(u\in \mathbf{A}[\mathbf{X}]\) , take \(g_{i}\in \mathbf{A}[\mathbf{Y}]\) for \(i\in I\) and put \(\mathbf{g} = (g_i)_{i\in I}\) ; let \(u(\mathbf{g})\in \mathbf{A}[\mathbf{Y}]\) be the polynomial obtained by substituting the polynomials \(g_{i}\) for \(\mathbf{X}_i\) in \(u\) . Let \(\mathbf{y} = (y_j)_{j\in J}\) be a family of pairwise permutable elements of a unital associative A-algebra \(\mathbf{E}'\) ; applying (1) and taking for \(\lambda\) the homomorphism \(g\mapsto g(\mathbf{y})\) of E into \(\mathbf{E}'\) , we obtain

\[
(u (\mathbf {g})) (\mathbf {y}) = u \left(\left(g _ {i} (\mathbf {y})\right)\right).\tag{2}
\]

If \(\mathbf{f} = (f_i)_{i\in I}\in (\mathbf{A}[(\mathbf{X}_j)_{j\in J}])^I\) and \(\mathbf{g} = (g_j)_{j\in J}\in (\mathbf{A}[(\mathbf{Y}_k)_{k\in K}])^J\) , we denote by \(\mathbf{f}\circ \mathbf{g}\) or \(f(\mathbf{g})\) , the family of polynomials \((f_i(\mathbf{g}))_{i\in I}\in (\mathbf{A}[(\mathbf{Y}_k)_{k\in K}])^I\) . If we denote by \(\tilde{\mathbf{f}}\) the mapping \(\mathbf{x} \mapsto (f_i(\mathbf{x}))_{i \in I}\) of \(\mathrm{E}^{\prime J}\) into \(\mathrm{E}^{\prime I}\) (where \(\mathrm{E}'\) is a unital associative and commutative A-algebra), then the relation (2) implies

\[
(\mathbf {f} \circ \mathbf {g}) ^ {\sim} = \tilde {\mathbf {f}} \circ \tilde {\mathbf {g}}.\tag{3}
\]

If \(\mathbf{h} = (h_k)_{k\in K}\in (\mathbf{A}[(\mathbf{Z}_l)_{l\in L}])^{\mathbf{K}}\) , it follows from (2) that :

\[
\mathbf {f} \circ (\mathbf {g} \circ \mathbf {h}) = (\mathbf {f} \circ \mathbf {g}) \circ \mathbf {h}.\tag{4}
\]

PROPOSITION 2. --- Let \(\mathbf{a} = (a_i)_{i \in I}\) be a family of elements of A and let \(u \in A[X]\) . If v is the polynomial obtained by substituting \(X_i + a_i\) for \(X_i\) for each \(i \in I\) , then the constant term of v is equal to \(u(\mathbf{a})\) .

The constant term of v is obtained by substituting 0 for \(X_{i}\) in v for each \(i \in I\) . The result therefore follows by (2).

COROLLARY 1. --- Let m be the ideal of polynomials \(u \in A[X]\) such that \(u(\mathbf{a}) = 0\) . Then m is generated by the polynomials \(X_{i} - a_{i}\) (for \(i \in I\) ).

It is clear that \(X_{i} - a_{i} \in m\) for each \(i \in I\) . Let \(u \in m\) and let v be as in Prop. 2. Since v has no constant term, there exists a family \((\mathbf{P}_{i})_{i \in I}\) of polynomials in A{[}X{]} with finite support such that

\[
v (\mathbf {X}) = \sum_ {i \in I} \mathrm{X} _ {i} \cdot \mathrm{P} _ {i} (\mathbf {X}).
\]

If we replace \(X_{i}\) by \(X_{i}-a_{i}\) for each \(i\in I\) in the above equation, we obtain a relation of the form \(u(\mathbf{X})=\sum_{i\in I}\left(\mathbf{X}_{i}-a_{i}\right)\cdot\mathbf{P}_{i}^{\prime}(\mathbf{X})\) , whence the corollary.

COROLLARY 2. --- Let \(\mathbf{X} = (\mathbf{X}_i)_{i \in I}\) and \(\mathbf{Y} = (\mathbf{Y}_i)_{i \in I}\) be two families of indeterminates. The set of polynomials \(u \in \mathbf{A}[\mathbf{X}, \mathbf{Y}]\) such that \(u(\mathbf{X}, \mathbf{X}) = 0\) is the ideal of \(\mathbf{A}[\mathbf{X}, \mathbf{Y}]\) generated by the polynomials \(\mathbf{X}_i - \mathbf{Y}_i\) (for \(i \in I\) ).

This corollary results directly from Cor. 1 on replacing A by A{[}Y{]} and \(a_{i}\) by \(Y_{i}\) , interpreting A{[}X, Y{]} as polynomial ring in the \(X_{i}\) with coefficients in A{[}Y{]}.

PROPOSITION 3. --- Let \(u \in \mathbf{A}[\mathbf{X}]\) and let \(\mathbf{X} \cdot \mathbf{Z}\) be the family \((\mathbf{X}_i\mathbf{Z})_{i \in I}\) of elements of the polynomial ring \(\mathbf{A}[\mathbf{X}][\mathbf{Z}]\) . The coefficient of \(\mathbf{Z}^k\) in \(u(\mathbf{X}, \mathbf{Z})\) is the homogeneous component of degree \(k\) of \(u\) , for every positive integer \(k\) .

It suffices to prove this Prop. in the case where \(u\) is a monomial, and in this case the result is clear.

COROLLARY. --- For a polynomial \(u \in A[X]\) to be homogeneous of degree \(k\) it is necessary and sufficient that :

\[
u (\mathbf {X} \cdot \mathbf {Z}) = u (\mathbf {X}) \cdot \mathbf {Z} ^ {k}.
\]

Remark. --- Let \(x \in A^{I}\) and let f be the mapping \(u \mapsto u(x)\) of A{[}X{]} into A. Given an A-module M, we consider the homomorphism \(1 \otimes f\) of M{[}X{]} = M \(\otimes_{A} A[X]\) into \(\mathbf{M} \otimes_{\mathbf{A}} \mathbf{A} = \mathbf{M}\) . For each \(v \in \mathbf{M}[\mathbf{X}]\) we have \((1 \otimes f)(v) = v(\mathbf{x})\) . If \(v = \sum_{\nu \in \mathbf{N}^{(I)}} e_{\nu} X^{\nu}\) , then \(v(\mathbf{x}) = \sum_{\nu \in \mathbf{N}^{(I)}} \mathbf{x}^{\nu} e_{\nu}\) .

\subsection{Differentials and derivations}

Let \(\mathbf{B} = \mathbf{A}[(\mathbf{X}_i)_{i\in I}]\) , then by III, p.~569 there exists for each \(i\in I\) one and only one A-derivation \(\mathbf{D}_i\) of \(\mathbf{B}\) such that

\[
\mathrm{D} _ {i} \mathrm{X} _ {i} = 1, \quad \mathrm{D} _ {i} \mathrm{X} _ {j} = 0 \quad \text { for } \quad j \neq i.\tag{5}
\]

The polynomial \(D_{i}P\) is called the partial derivative of P with respect to \(X_{i}\) ; we shall also denote it by \(D_{X_{i}}P\) or \(\frac{\partial P}{\partial X_{i}}\) or \(P'_{X_{i}}\) . By III, p.~558, formula (21), we have, for \(\nu = (\nu_{j}) \in \mathbf{N}^{(I)}\) ,

\[
\mathrm{D} _ {i} \left(\mathbf {X} ^ {\nu}\right) = \left\{ \begin{array}{c c c} \nu_ {i} \mathbf {X} _ {i} ^ {\nu_ {i} - 1} \prod_ {j \in I - \{i \}} \mathbf {X} _ {j} ^ {\nu_ {j}} & \text {if} & \nu_ {i} > 0 \\ 0 & \text {if} & \nu_ {i} = 0. \end{array} \right.\tag{6}
\]

It follows from (6) that \(\mathbf{D}_i\mathbf{D}_j = \mathbf{D}_j\mathbf{D}_i\) for any \(i, j \in I\) . For \(\nu = (\nu_i)_{i \in I} \in \mathbf{N}^{(I)}\) we shall put \(\mathbf{D}^{\nu} = \prod_{i \in I} \mathbf{D}_i^{\nu_i}\) and \(\nu! = \prod_{i \in I} (\nu_i!)\) . With the product ordering on \(\mathbf{N}^{(I)}\) we have

\[
\mathrm{D} ^ {\nu} \left(\mathrm{X} ^ {\mu}\right) = \left\{ \begin{array}{c l} \frac {\mu !}{(\mu - \nu) !} \mathrm{X} ^ {\mu - \nu} & \text { if } \quad \nu \leqslant \mu , \\ 0 & \text { if   not }. \end{array} \right.
\]

When P is a polynomial in a single indeterminate X, the unique partial derivative of P is written DP or \(\frac{dP}{dX}\) or \(P'\) and is called simply the derivative of P.

Again let \(\mathbf{B} = \mathbf{A}\left[(\mathbf{X}_i)_{i\in I}\right]\) ; by III, p.~569 the B-module of A-differentials of B, \(\Omega_{\mathbf{A}}(\mathbf{B})\) , has the family \((d\mathbf{X}_i)_{i\in I}\) of differentials of the \(\mathbf{X}_i\) as basis. Let \(\partial_i\) be the coordinate form of index \(i\) relative to this basis over \(\Omega_{\mathbf{A}}(\mathbf{B})\) . Then the mapping \(u\mapsto \langle \partial_i,du\rangle\) of B into itself is a derivation of B which maps \(\mathbf{X}_i\) to 1 and \(\mathbf{X}_j\) to 0 for \(j\neq i\) , and hence is \(\mathbf{D}_i\) ; in other words, we have

\[
d u = \sum_ {i \in I} \left(\mathrm{D} _ {i} u\right) d \mathrm{X} _ {i}\tag{7}
\]

for each \(u \in \mathbf{B}\) . If \(\mathbf{I}\) is finite, \((\mathbf{D}_i)_{i \in \mathbf{I}}\) is a basis of the B-module of derivations of \(\mathbf{B}\) .

PROPOSITION 4. --- Let E be an associative, commutative and unital A-algebra, \(\mathbf{x} = (x_i)_{i\in I}\) a family of elements of E, \(u\) an element of \(\mathbf{A}[(\mathbf{X}_i)_{i\in I}]\) and \(y = u(\mathbf{x})\) . Then for every derivation D of E into an E-module we have

\[
\mathrm{D} y = \sum_ {i \in I} \left(\mathrm{D} _ {i} u\right) (\mathbf {x}) \cdot \mathrm{D} x _ {i}.
\]

It suffices to prove the proposition when \(u\) is a monomial, and in that case it follows from III, p.~558, Prop. 6.

COROLLARY. --- Let \(f \in \mathbf{A}[\mathbf{X}_1, \ldots, \mathbf{X}_p]\) and \(g_i \in \mathbf{A}[\mathbf{Y}_1, \ldots, \mathbf{Y}_q]\) for \(1 \leqslant i \leqslant p\) , and write \(h = f(g_1, \ldots, g_p)\) , then for \(1 \leqslant j \leqslant q\) we have

\[
\frac {\partial h}{\partial \mathbf {Y} _ {j}} = \sum_ {i = 1} ^ {p} \mathrm{D} _ {i} f (g _ {1}, \dots , g _ {p}) \cdot \frac {\partial g _ {i}}{\partial \mathbf {Y} _ {j}}.\tag{8}
\]

This is the special case \(\mathbf{E} = \mathbf{A}[\mathbf{Y}_1,\dots ,\mathbf{Y}_q]\) , \(x_{i} = g_{i}\) and \(\mathbf{D} = \partial /\partial \mathbf{Y}_j\) of Prop. 4.

Let \(\mathbf{X} = (\mathbf{X}_i)_{i\in I}\) , \(\mathbf{Y} = (\mathbf{Y}_i)_{i\in I}\) be two disjoint families of indeterminates, and write \(\mathbf{X} + \mathbf{Y}\) for the family \((\mathbf{X}_i + \mathbf{Y}_i)_{i\in I}\) . Given \(u\in \mathbf{A}[\mathbf{X}]\) , consider the element \(u(\mathbf{X} + \mathbf{Y})\) of \(\mathbf{A}[\mathbf{X},\mathbf{Y}]\) . For \(\nu \in \mathbf{N}^{(I)}\) we denote by \(\Delta^{\nu}u\) the coefficient of \(\mathbf{Y}^{\nu}\) in \(u(\mathbf{X} + \mathbf{Y})\) , considered as polynomial in the \(\mathbf{Y}_i\) with coefficients in \(\mathbf{A}[\mathbf{X}]\) . By definition we have \(\Delta^{\nu}u\in \mathbf{A}[\mathbf{X}]\) and

\[
u (\mathbf {X} + \mathbf {Y}) = \sum_ {\nu} \left(\Delta^ {\nu} u\right) (\mathbf {X}) \mathbf {Y} ^ {\nu}.\tag{9}
\]

(Here and in the rest of this No., the summations are extended over the index set \(\mathbf{N}^{(I)}\) unless the contrary is said.)

Let \(\mathbf{a} \in \mathbf{A}^{I}\) , then on substituting \(\mathbf{a}\) for \(\mathbf{X}\) and \(\mathbf{X} - \mathbf{a}\) for \(\mathbf{Y}\) in (9) we obtain

\[
u (\mathbf {X}) = \sum_ {\nu} \left(\Delta^ {\nu} u\right) (\mathbf {a}) (\mathbf {X} - \mathbf {a}) ^ {\nu}.\tag{10}
\]

In particular we have

\[
u (\mathbf {X}) = \sum_ {\nu} \left(\Delta^ {\nu} u\right) (0) \mathbf {X} ^ {\nu}.\tag{11}
\]

If \(u, v \in \mathbf{A}[\mathbf{X}]\) , we have

\[
\begin{array}{r l} (u v) (\mathbf {X} + \mathbf {Y}) & = \left(\sum_ {\nu} (\Delta^ {\nu} u) (\mathbf {X}) \mathbf {Y} ^ {\nu}\right) \left(\sum_ {\rho} (\Delta^ {\rho} v) (\mathbf {X}) \mathbf {Y} ^ {\rho}\right) \\ & = \sum_ {\sigma} \left[ \sum_ {\nu + \rho = \sigma} (\Delta^ {\nu} u) (\mathbf {X}) (\Delta^ {\rho} v) (\mathbf {X}) \right] \mathbf {Y} ^ {\sigma} \end{array}
\]

hence

\begin{enumerate}
\def\labelenumi{(\arabic{enumi})}
\setcounter{enumi}{11}
\tightlist
\item
\end{enumerate}

\[
\Delta^ {\sigma} (u v) = \sum_ {\nu + \rho = \sigma} (\Delta^ {\nu} u) (\Delta^ {\rho} v).
\]

Let \(\mathbf{Z} = (Z_i)_{i\in I}\) be another family of indeterminates.

We have :

\[
\begin{array}{r l} \sum_ {\nu} (\Delta^ {\nu} u) (\mathbf {X}) (\mathbf {Y} + \mathbf {Z}) ^ {\nu} & = u (\mathbf {X} + \mathbf {Y} + \mathbf {Z}) \\ & = \sum_ {\sigma} (\Delta^ {\sigma} u) (\mathbf {X} + \mathbf {Y}) \mathbf {Z} ^ {\sigma} \\ & = \sum_ {\rho , \sigma} (\Delta^ {\rho} \Delta^ {\sigma} u) (\mathbf {X}) \mathbf {Y} ^ {\rho} \mathbf {Z} ^ {\sigma}, \end{array}
\]

hence by I, p.~99, Cor. 2 :

\[
\Delta^ {\rho} \Delta^ {\sigma} u = \frac {(\rho + \sigma) !}{\rho ! \sigma !} \Delta^ {\rho + \sigma} u.\tag{13}
\]

PROPOSITION 5. --- For any \(u \in \mathbf{A}[\mathbf{X}]\) and \(\nu \in \mathbf{N}^{(I)}\) we have

\[
\mathbf {D} ^ {\nu} \boldsymbol {u} = \nu ! \Delta^ {\nu} \boldsymbol {u}.
\]

Suppose first that \(\nu\) has length 1; then there exists an element \(i\) of I such that \(\nu = \varepsilon_i\) , that is, \(\nu_i = 1\) and \(\nu_j = 0\) for all \(j \neq i\) in I. Formula (12) shows that \(\Delta^{\varepsilon_i}\) is a derivation of the A-algebra A{[}X{]} which clearly has the value zero on \(\mathbf{X}_j\) for \(j \neq i\) and the value 1 on \(\mathbf{X}_i\) . We thus have \(\Delta^{\varepsilon_i} = \mathbf{D}_i\) for each \(i \in I\) .

By (13) we have

\[
(\rho ! \Delta^ {\rho}) \cdot (\sigma ! \Delta^ {\sigma}) = (\rho + \sigma)! \Delta^ {\rho + \sigma}\tag{14}
\]

in the endomorphism algebra of the A-module \(\mathbf{A}[\mathbf{X}]\) . It now follows by induction on the length of \(\nu\) that \(\nu! \Delta^{\nu} = \mathbf{D}^{\nu}\) .

If \(\mathbf{A}\) is a \(\mathbf{Q}\) -algebra, the formulae (9), (10), (11) may thus be written

\begin{enumerate}
\def\labelenumi{(\arabic{enumi})}
\setcounter{enumi}{14}
\tightlist
\item
\end{enumerate}

\[
u (\mathbf {X} + \mathbf {Y}) = \sum_ {\nu} \frac {1}{\nu !} (\mathrm{D} ^ {\nu} u) (\mathbf {X}) \mathbf {Y} ^ {\nu}
\]

\[
u (\mathbf {X}) = \sum_ {\nu} \frac {1}{\nu !} (D ^ {\nu} u) (\mathbf {a}) (\mathbf {X} - \mathbf {a}) ^ {\nu}\tag{16}
\]

\[
u (\mathbf {X}) = \sum_ {\nu} \frac {1}{\nu !} (D ^ {\nu} u) (0) \mathbf {X} ^ {\nu}.\tag{17}
\]

The formulae (15), (16), (17) are all three called « Taylor's formula ».

PROPOSITION 6 (« Euler's identity »). --- Let \(u \in \mathbf{A}[\mathbf{X}]\) be a homogeneous polynomial of degree \(r\) ; we have

\[
\sum_ {i \in I} X _ {i} \cdot D _ {i} u = r u.
\]

Let D be the A-linear mapping of A{[}X{]} into itself such that \(D(v) = sv\) when v is homogeneous of degree s. We know (III, p.~554, Example 6) that D is a derivation of A{[}X{]}. Thus Prop. 6 is a Corollary of Prop. 4 (IV, p.~6).

\subsection{Divisors of zero in a polynomial ring}

PROPOSITION 7. --- Let \(f \in A[X]\) be a non-zero polynomial in an indeterminate and \(\alpha\) its leading coefficient. If \(\alpha\) is cancellable in \(A\) (in particular if \(f\) is monic) then we have for any non-zero element \(g\) of \(A[X]\) ,

\[
f g \neq 0 \quad a n d \quad \deg (f g) = \deg f + \deg g.
\]

Let \(g \in A[X]\) be a non-zero polynomial, \(\beta\) its leading coefficient, \(n = \deg f\) and \(p = \deg g\) . Then the coefficient of \(X^{n+p}\) in \(fg\) is \(\alpha\beta\) which does not vanish, whence the proposition.

PROPOSITION 8. -- If A is an integral domain, then so is \(\mathbf{A}[(\mathbf{X}_i)_{i\in I}]\) .

Let \(u, v\) be two non-zero elements of \(\mathbf{A}[(\mathbf{X}_i)_{i \in I}]\) ; we have to show that \(uv \neq 0\) . Now \(u\) and \(v\) belong to a ring \(\mathbf{A}[(\mathbf{X}_j)_{j \in J}]\) where \(J\) is a finite subset of \(I\) . Thus we can limit ourselves to the case when \(I\) is finite and equal to \(\{1, 2, \ldots, p\}\) . On the other hand, the ring \(\mathbf{A}[\mathbf{X}_1, \ldots, \mathbf{X}_p]\) is isomorphic to the polynomial ring in \(\mathbf{X}_p\) with coefficients in \(\mathbf{A}[\mathbf{X}_1, \ldots, \mathbf{X}_{p-1}]\) . By induction on \(p\) we are thus reduced to proving the proposition for \(\mathbf{A}[\mathbf{X}]\) , and now it suffices to apply Prop. 7.

COROLLARY 1. --- If A is an integral domain, and u, v are elements of \(\mathbf{A}[(\mathbf{X}_{i})_{i\in I}]\) , then \(\deg(uv)=\deg u+\deg v\) .

We can limit ourselves to the case where \(u\) and \(v\) are non-zero. Let \(m = \deg u\) , \(n = \deg v\) ; we have

\[
u = u _ {0} + u _ {1} + \dots + u _ {m}, v = v _ {0} + v _ {1} + \dots + v _ {n}
\]

where \(u_{h}\) (resp. \(v_{h}\) ) is the homogeneous component of degree h of u (resp. v). Since \(u_{m} \neq 0\) and \(v_{n} \neq 0\) , we have \(u_{m}v_{n} \neq 0\) (Prop. 8). Now \(uv = u_{m}v_{n} + w\) with \(\deg w < m + n\) , whence the result.

COROLLARY 2. --- If A is an integral domain, the invertible elements of A{[}(X \(_{i}\) ) \(_{i \in I}\) {]} are the invertible elements of A.

This follows immediately from Cor. 1.

PROPOSITION 9. --- Let \(u \in A[(X_i)_{i \in I}]\) ; then for \(u\) to be nilpotent in the ring \(A[(X_i)_{i \in I}]\) it is necessary and sufficient for all its coefficients to be nilpotent in the ring \(A\) .

As in the proof of Prop. 8 we can make a reduction to the case of polynomials in one variable X. If all the coefficients of u are nilpotent, then u is nilpotent (I, p.~99, Cor. 1). Suppose that \(u\) is nilpotent but not zero and let \(n\) be its degree; we shall argue by induction on \(n\) . Let \(\alpha\) be the leading coefficient of \(u\) . There exists an integer \(m > 0\) such that \(u^m = 0\) . The leading coefficient of \(u^m\) is \(\alpha^m\) , hence \(\alpha^m = 0\) . Now \(u - \alpha X^n\) is nilpotent (I, loc. cit.) and the induction hypothesis shows all the coefficients of \(u - \alpha X^n\) to be nilpotent. Thus all the coefficients of \(u\) are nilpotent.

Remark. --- Let u and v be elements of \(\mathbf{A}[(\mathbf{X}_{i})_{i\in I}]\) , and suppose that A is an integral domain, v is a non-zero multiple of u and v is homogeneous; then u is also homogeneous. For let \(u'\in\mathbf{A}[(\mathbf{X}_{i})_{i\in I}]\) be such that \(v=uu'\) ; we have \(u\neq0\) , \(u'\neq0\) , and if

\[
\begin{array}{c} u = u _ {h} + u _ {h + 1} + \dots + u _ {k} \\ u ^ {\prime} = u _ {h ^ {\prime}} ^ {\prime} + u _ {h ^ {\prime} + 1} ^ {\prime} + \dots + u _ {k ^ {\prime}} ^ {\prime} \end{array}
\]

are the decompositions of u and \(u'\) into homogeneous components, with \(u_{h} \neq 0\) , \(u_{k} \neq 0\) , \(u_{h'}' \neq 0\) , \(u_{k'}' \neq 0\) , then \(v = u_{h} u_{h'}' + u_{h} u_{h'+1}' + \cdots + u_{k} u_{k'}'\) and \(u_{h} u_{h'}'\) is non-zero homogeneous of degree \(h + h'\) while \(u_{k} u_{k'}'\) is non-zero homogeneous of degree \(k + k'\) (Prop. 8). Since v is homogeneous, we have \(h + h' = k + k'\) whence h = k, \(h' = k'\) .

\subsection{Euclidean division of polynomials in one indeterminate}

PROPOSITION 10. --- Let \(f\) and \(g\) be non-zero elements of \(\mathbf{A}[\mathbf{X}]\) of degrees \(m\) and \(n\) respectively. Let \(\alpha_0\) be the leading coefficient of \(f\) and \(\mu = \sup (n - m + 1, 0)\) . There exist \(u, v \in \mathbf{A}[\mathbf{X}]\) such that

\[
\alpha_ {0} ^ {\mu} g = u f + v, \quad \deg v <   m.
\]

If \(\alpha_0\) is cancellable in A, then \(u\) and \(v\) are uniquely determined by these properties.

The existence of \(u\) and \(v\) is clear when \(n < m\) , because then we can take \(u = 0\) and \(v = g\) . For \(n \geqslant m\) we shall use induction on \(n\) . Let \(\beta\) be the leading coefficient of \(g\) ; if \(f = \sum_{k=0}^{m} \alpha_k X^{m-k}\) , we can write \(\alpha_0^\mu g = \alpha_0^{\mu-1} \beta X^{n-m} f + \alpha_0^{\mu-1} g_1\) , where \(g_1 \in A[X]\) and \(\deg g_1 < n\) . By the induction hypothesis there exist \(u_1, v \in A[X]\) such that \(\alpha_0^{\mu-1} g_1 = u_1 f + v\) and \(\deg v < m\) . Hence

\[
\alpha_ {0} ^ {\mu} g = \left(\alpha_ {0} ^ {\mu - 1} \beta X ^ {n - m} + u _ {1}\right) f + v
\]

and it suffices to put \(u = \alpha_0^{\mu - 1}\beta X^{n - m} + u_1\) .

Assuming \(\alpha_0\) to be cancellable in \(\mathbf{A}\) , let us now prove the uniqueness of \(u\) and \(v\) . Let \(u, v, u_1, v_1 \in \mathbf{A}[\mathbf{X}]\) be such that

\[
\alpha_ {0} ^ {\mu} g = u f + v = u _ {1} f + v _ {1}, \quad \deg v <   m, \quad \deg v _ {1} <   m.
\]

We have \((u - u_{1})f = v_{1} - v\) and \(\deg (v_1 - v) < m\) , hence \(u - u_{1} = 0\) (IV, p.~9, Prop. 7) and therefore \(v_{1} - v = 0\) .

COROLLARY (« Euclidean division of polynomials »). --- Let \(f\) be a non-zero element of A{[}X{]} whose leading coefficient is invertible and \(m = \deg f\) .

\begin{enumerate}
\def\labelenumi{(\roman{enumi})}
\tightlist
\item
  For every \(g \in \mathbf{A}[\mathbf{X}]\) there exist \(u, v \in \mathbf{A}[\mathbf{X}]\) such that
\end{enumerate}

\[
g = u f + v, \quad \deg v <   m.
\]

Moreover, these conditions determine \(u\) and \(v\) uniquely.

\begin{enumerate}
\def\labelenumi{(\roman{enumi})}
\setcounter{enumi}{1}
\item
  The sub-A-modules \(A + AX + \cdots + AX^{m-1}\) and \(fA[X]\) of A{[}X{]} are supplementary in A{[}X{]}.
\item
  Assume \(f\) non-constant and consider \(\mathbf{A}[\mathbf{X}]\) as an \(\mathbf{A}[\mathbf{T}]\) -module by means of the homomorphism \(u(\mathbf{T}) \mapsto u(f(\mathbf{X}))\) of \(\mathbf{A}[\mathbf{T}]\) into \(\mathbf{A}[\mathbf{X}]\) . Then \(\mathbf{A}[\mathbf{X}]\) is a free \(\mathbf{A}[\mathbf{T}]\) -module with basis \((1, \mathbf{X}, \ldots, \mathbf{X}^{m-1})\) .
\end{enumerate}

Assertions (i) and (ii) are immediate consequences of Prop. 10.

Let us prove (iii). Let \(\psi\) be the homomorphism \(v \mapsto v(f(X), X)\) of A{[}T, X{]} into A{[}X{]}. Consider A{[}T, X{]} first as polynomial ring in T with coefficients in A{[}X{]}; Cor. 1 of IV, p.~5 shows that the kernel \(\mathfrak{a}\) of \(\psi\) is the ideal (T - f(X)) of A{[}T, X{]}. Consider now A{[}T, X{]} as polynomial ring in X with coefficients in A{[}T{]}; then \(\psi\) is an A{[}T{]}-linear mapping of A{[}T{]}{[}X{]} into A{[}X{]}. Assertion (ii) above (applied to the polynomial \(f(X) - T\) in X with coefficients in A{[}T{]}) shows that \((1, X, ..., X^{m-1})\) is a basis of an A{[}T{]}-submodule of A{[}T, X{]} supplementary to \(\mathfrak{a}\) . Since \(\psi(X^i) = X^i\) for every integer \(i \geqslant 0\) , (iii) follows at once.

With the notations of (i) we shall say that u is the quotient and v the remainder in the Euclidean division of g by f; for the remainder to vanish it is necessary and sufficient that f should divide g.

\subsection{\texorpdfstring{Divisibility of polynomials in one indeterminate \(^{1}\)}{Divisibility of polynomials in one indeterminate [1]}}

PROPOSITION 11. --- Let K be a commutative field.

\begin{enumerate}
\def\labelenumi{(\roman{enumi})}
\item
  For every non-zero ideal \(\mathfrak{a}\) of \(\mathbf{K}[\mathbf{X}]\) there exists precisely one monic polynomial \(f\) in \(\mathbf{K}[\mathbf{X}]\) such that \(\mathfrak{a} = (f)\) .
\item
  Let \(f_{1}\) and \(f_{2}\) be in \(\mathbf{K}[\mathbf{X}]\) ; for \((f_{1}) = (f_{2})\) to hold it is necessary and sufficient that there exist a non-zero element \(\lambda\) of \(\mathbf{K}\) such that \(f_{2} = \lambda f_{1}\) .
\end{enumerate}

Let us prove (ii), the sufficiency of the stated condition being clear. The case where \(f_{1}\) and \(f_{2}\) generate the zero ideal is trivial. Thus assume that the non-zero polynomials \(f_{1}\) and \(f_{2}\) generate the same ideal of K{[}X{]}. Then there exist polynomials \(u_{1}\) and \(u_{2}\) such that \(f_{1}=u_{1}f_{2}\) and \(f_{2}=u_{2}f_{1}\) ; it follows that \(u_{1}u_{2} = 1\) , whence \(\deg u_{1} + \deg u_{2} = 0\) and so \(\deg u_{2} = 0\) . We have thus shown that \(u_{2}\) is a non-zero element of \(\mathbf{K}\) .

To prove (i), let \(f\) be a monic polynomial in \(\mathfrak{a}\) of least possible degree. Given \(g\) in \(\mathfrak{a}\) , let \(u\) and \(v\) be the quotient and remainder of the Euclidean division of \(g\) by \(f\) ; then \(v = g - uf\) belongs to \(\mathfrak{a}\) and we have \(\deg v < \deg f\) ; if \(v\) were non-zero, there would be a non-zero element \(\lambda\) of \(K\) such that \(\lambda v\) is monic, and since \(\lambda v \in \mathfrak{a}\) , this would contradict the definition of \(f\) . We thus have \(\mathfrak{a} = (f)\) ; the uniqueness of the monic polynomial \(f\) such that \(\mathfrak{a} = (f)\) now follows from (ii).

PROPOSITION 12. --- Let K be a commutative field and f, g two elements of K{[}X{]}. For every polynomial d in K{[}X{]} the following properties are equivalent:

\begin{enumerate}
\def\labelenumi{(\roman{enumi})}
\item
  The polynomial \(d\) divides \(f\) and \(g\) and every polynomial which divides both \(f\) and \(g\) divides \(d\) .
\item
  The polynomial \(d\) divides \(f\) and \(g\) and there exist two polynomials \(u\) and \(v\) such that \(d = uf + vg\) .
\item
  The relation \((d) = (f) + (g)\) holds between ideals in \(\mathbf{K}[\mathbf{X}]\) .
\end{enumerate}

The polynomial d is determined up to multiplication by a non-zero element of K by these properties. If f and g are not both zero, then \(d \neq 0\) and the degree of d majorizes the degree of every polynomial dividing both f and g.

When \(f\) and \(g\) are zero, each of the properties (i) to (iii) is satisfied only for \(d = 0\) , hence they are then equivalent. Henceforth we assume that \(f, g\) are not both 0 and we denote by \(\mathfrak{a}\) the ideal \((f) + (g)\) of \(\mathbf{K}[\mathbf{X}]\) .

We remark that for any polynomials \(u\) and \(v\) in \(\mathbf{K}[\mathbf{X}]\) the properties \((u) \supset (v)\) and \(\langle u \text{ divides } v \rangle\) are equivalent. The assertion (ii) is thus equivalent to \(\langle (d) \supset (f) \text{ and } (d) \supset (g) \text{ and } d \in (f) + (g) \rangle\) , that is (iii). It is clear that (ii) implies (i). Finally suppose that (i) holds; we have \((d) \supset (f)\) and \((d) \supset (g)\) , whence \((d) \supset \mathfrak{a}\) ; on the other hand, by Prop. 11 (IV, p.~12) there exists a polynomial \(d_1\) such that \(\mathfrak{a} = (d_1)\) ; since \(d_1\) divides both \(f\) and \(g\) , it divides \(d\) by hypothesis, whence \((d) \subset \mathfrak{a}\) , and finally we have \((d) = \mathfrak{a}\) , that is, (iii).

The other assertions of Prop. 12 are immediate consequences of Prop. 11 applied to the ideal \(\mathfrak{a} = (f) + (g)\) .

DEFINITION 1. --- With the notation of Prop. 12 we say that d is a greatest common divisor (gcd for short) of f and g. We say that f and g are relatively prime or that f is prime to g if 1 is a gcd of f and g.

To say that f and g are relatively prime thus means that there exist polynomials u and v in K{[}X{]} such that \(uf + vg = 1\) .

COROLLARY 1. --- Let \(d\) be a gcd of \(f\) and \(g\) and \(\mathbf{K}'\) a commutative field containing \(\mathbf{K}\) as subfield. Then \(d\) is a gcd of \(f\) and \(g\) considered as elements of \(\mathbf{K}'[\mathbf{X}]\) .

This follows from Prop. 12, (iii).

COROLLARY 2. --- Let d be a gcd of f and g.

\begin{enumerate}
\def\labelenumi{(\roman{enumi})}
\item
  If \(u \in \mathbf{K}[\mathbf{X}]\) , du is a gcd of fu and gu.
\item
  If \(v \in \mathbf{K}[\mathbf{X}]\) is a divisor ( \(\neq 0\) ) of \(f\) and \(g\) , then \(d / v\) is a gcd of \(f / v\) and \(g / v\) .
\end{enumerate}

This follows from Prop. 12, (ii).

COROLLARY 3. --- Let w be a common factor of f and g. For w to be a gcd of f and g it is necessary and sufficient that f/w and g/w are relatively prime. This follows from Cor. 2.

COROLLARY 4. --- Let \(f, g, h \in \mathbf{K}[\mathbf{X}]\) . If \(f\) divides \(gh\) and is prime to \(g\) , then \(f\) divides \(h\) .

For \(f\) divides \(gh\) and \(fh\) , hence \(f\) divides every gcd of \(gh\) and \(fh\) , in particular \(h\) (Cor. 2, (i)).

COROLLARY 5. --- Let \(f, g \in \mathbf{K}[\mathbf{X}]\) . For \(f\) and \(g\) to be relatively prime it is necessary and sufficient that the canonical image of \(g\) in \(\mathbf{K}[\mathbf{X}] / (f)\) should be invertible.

For this condition means that there exist \(u, v \in \mathbf{K}[\mathbf{X}]\) such that \(uf + vg = 1\) .

COROLLARY 6. --- Let \(f, g_1, g_2, \ldots, g_n \in \mathbf{K}[\mathbf{X}]\) . If \(f\) is prime to \(g_1, g_2, \ldots, g_n\) , then \(f\) is prime to \(g_1g_2\ldots g_n\) .

* COROLLARY 7. --- For \(f\) and \(g\) to be relatively prime it is necessary and sufficient that they have no common roots in any extension of \(K\) .

For if \(d\) is a gcd of \(f, g\) then the roots common to \(f\) and \(g\) in an extension \(\mathbf{K}'\) of \(\mathbf{K}\) are the roots of \(d\) in \(\mathbf{K}'\) . Now the corollary follows from \(V\) , p.~21, Prop. 4. *

\subsection{Irreducible polynomials}

DEFINITION 2. --- Let K be a commutative field. We say that \(f \in \mathbf{K}[\mathbf{X}]\) is irreducible if \(\deg f \geqslant 1\) and \(f\) is not divisible by any polynomial \(g\) such that \(0 < \deg g < \deg f\) .

It comes to the same to say that \(\deg f \geqslant 1\) and the only divisors of f in K{[}X{]} are the scalars \(\neq 0\) and the products of f by scalars \(\neq 0\) . Since the relation \((f) \subset (g)\) means that g divides f, we see that the irreducible polynomials of K{[}X{]} may also be defined as the polynomials f such that the ideal \((f)\) is maximal (I, p.~104).

Let \(f, g \in \mathbf{K}[\mathbf{X}]\) . If \(f\) is irreducible, it is clear that either \(f\) and \(g\) are relatively prime or that \(f\) divides \(g\) . If \(f\) and \(g\) are irreducible, either \(f\) and \(g\) are relatively prime or each is a product of the other by a scalar \(\neq 0\) . In particular, two distinct irreducible monic polynomials are relatively prime.

PROPOSITION 13. --- Let \(\mathcal{I}\) be the set of irreducible monic polynomials in K{[}X{]}. Let f be a non-zero element of K{[}X{]} and α its leading coefficient; then there exists precisely one family of positive integers \((v_{p})_{p\in \mathcal{I}}\) with finite support, such that we have a decomposition

\[
f = \alpha \prod_ {p \in \mathcal {I}} p ^ {v _ {p}}.\tag{18}
\]

It suffices to prove the proposition when \(f\) is monic, that is when \(\alpha = 1\) . We shall argue by induction on the degree \(n\) of \(f\) , the case \(n = 0\) being trivial. Suppose then that \(n \geqslant 1\) and that the proposition has been established for all polynomials of degree \(< n\) .

Let E be the set of monic polynomials \(\neq1\) which divide f; we have \(f\in E\) hence E is not empty and there exists in E a polynomial g of least degree. It is clear that g is irreducible and there exists a monic polynomial h of degree \textless n such that f=gh; by the induction hypothesis h is the product of a finite family of irreducible monic polynomials, hence f has the same property. This proves the existence of the decomposition (18).

Let us now prove the uniqueness of the decomposition (18). Let \((w_{p})_{p\in \mathcal{I}}\) be a family of positive integers with finite support, such that \(f = \prod_{p\in \mathcal{I}}p^{w_p}\) . Since \(f\) is of

degree \(n \geqslant 1\) , there exists \(p \in \mathcal{I}\) such that \(w_{p} > 0\) ; if we had \(v_{p} = 0\) , f would be the product of a family of elements of \(\mathcal{I}\) distinct from p, hence it would be prime to p (IV, p.~13, Cor. 6), contrary to the fact that p divides f.~By the induction hypothesis the polynomial f/p admits a unique decomposition of type (18); hence we conclude the equality \(w_{q} = v_{q}\) for every \(q \in \mathcal{I}\) .

Let \(f\) be a non-zero polynomial in \(\mathbf{K}[\mathbf{X}]\) . We shall say that \(f\) has no multiple factors if the exponents \(v_{p}\) in the decomposition (18) are all \(\leqslant 1\) ; it comes to the same to say that \(f\) is the product of a finite sequence of pairwise distinct irreducible polynomials, or also that \(f\) is not divisible by the square of any non-constant polynomial of \(\mathbf{K}[\mathbf{X}]\) .

\section{ZEROS OF POLYNOMIALS}

\subsection{Roots of a polynomial in one indeterminate. Multiplicity}

Let \(g \in A[(X_i)_{i \in I}]\) and let \(E\) be a unital associative \(A\) -algebra. Let \(x = (x_i)_{i \in I}\) be a family of pairwise permutable elements of \(E\) . We shall say that \(x\) is a zero of \(g\) in \(E^{I}\) if \(g(x) = 0\) . If \(f\) is a polynomial in a single indeterminate, a zero of \(f\) in \(E\) is also called a root of \(f\) in \(E\) .

PROPOSITION 1. --- Let \(f \in \mathbf{A}[\mathbf{X}]\) and \(\alpha \in \mathbf{A}\) . The remainder of the division of \(f\) by \(\mathbf{X} - \alpha\) is \(f(\alpha)\) . For \(\alpha\) to be a root of \(f\) it is necessary and sufficient that \(\mathbf{X} - \alpha\) should be a divisor of \(f\) in \(\mathbf{A}[\mathbf{X}]\) .

For if \(u, v \in \mathbf{A}[\mathbf{X}]\) are such that \(f = (\mathbf{X} - \alpha)u + v\) , \(\deg v < 1\) , then \(v\) is a scalar and \(f(\alpha) = (\alpha - \alpha)u(\alpha) + v = v\) . This proves the first assertion and the second follows from it.

PROPOSITION 2. --- Let \(f \in \mathbf{A}[\mathbf{X}]\) , \(\alpha \in \mathbf{A}\) , and let \(h\) be an integer \(\geqslant 0\) . The following conditions are equivalent:

\begin{enumerate}
\def\labelenumi{(\roman{enumi})}
\item
  \(f\) is divisible by \((\mathbf{X} - \alpha)^h\) but not by \((\mathbf{X} - \alpha)^{h + 1}\) ;
\item
  there exists \(g \in \mathbf{A}[\mathbf{X}]\) such that \(f = (\mathbf{X} - \alpha)^h g\) and \(g(\alpha) \neq 0\) .
\item
  \(\Rightarrow\) (ii) follows at once from Prop. 1.
\item
  \(\Rightarrow\) (i): Suppose that \(f = (\mathbf{X} - \alpha)^h g\) , where \(g\) does not admit \(\alpha\) as a root. Then \(f\) is divisible by \((\mathbf{X} - \alpha)^h\) ; if \(g_1 \in \mathbf{A}[\mathbf{X}]\) existed such that \(f = (\mathbf{X} - \alpha)^{h + 1}g_1\) , then since \((\mathbf{X} - \alpha)^n\) is not a divisor of zero in \(\mathbf{A}[\mathbf{X}]\) (IV, p.~9, Prop. 7), we have \(g = (\mathbf{X} - \alpha)g_1\) and so \(g(\alpha) = 0\) , which is absurd.
\end{enumerate}

PROPOSITION 3. --- Let \(f\) be a non-zero element of \(\mathbf{A}[\mathbf{X}]\) and \(\alpha \in \mathbf{A}\) . There exists just one integer \(h \geqslant 0\) satisfying the conditions (i) and (ii) of Prop. 2.

This is clear for condition (i), bearing in mind the fact that if \(f\) is divisible by \((\mathbf{X} - \alpha)^h\) , then \(\deg f \geqslant h\) (IV, p.~9, Prop. 7).

DEFINITION 1. --- With the above notation we say that \(\alpha\) is of order \(h\) , or multiplicity \(h\) relative to \(f\) .

If h \textgreater{} 0 we also say that \(\alpha\) is a root of order h or multiplicity h of f.~A root of order 1 is called a simple root, a root of order 2 a double root,\ldots{} A root of order \textgreater{} 1 is said to be multiple.

Remarks. --- 1) If \(f=0\) we agree to say that \(\alpha\) has order \(\geqslant h\) relative to \(f\) , whatever \(\alpha\in\mathbf{A}\) and the integer \(h\geqslant0\) . For any \(f\in\mathbf{A}[\mathbf{X}]\) and \(\alpha\in\mathbf{A}\) , to say that \(\alpha\) has order \(\geqslant h\) relative to \(f\) means that \((\mathbf{X}-\alpha)^{h}\) divides \(f\) .

\begin{enumerate}
\def\labelenumi{\arabic{enumi})}
\setcounter{enumi}{1}
\tightlist
\item
  Let B be a commutative ring containing A as subring. Let \(f \in A[X]\) be nonzero and \(\alpha \in A\) . The order of \(\alpha\) relative to f is the same, whether we consider f as element of B{[}X{]} or as element of A{[}X{]}. This is clear from condition (ii) of Prop. 2.
\end{enumerate}

PROPOSITION 4. --- Let f and g be non-zero elements of A{[}X{]}. Let \(\alpha \in \mathbf{A}\) , and let the orders of \(\alpha\) relative to f and g be p and q respectively.

\begin{enumerate}
\def\labelenumi{(\roman{enumi})}
\item
  The order of \(\alpha\) relative to \(f + g\) is \(\geqslant \inf(p, q)\) . It is equal to \(\inf(p, q)\) if \(p \neq q\) .
\item
  The order of \(\alpha\) relative to \(fg\) is \(\geqslant p + q\) . It is equal to \(p + q\) if \(\mathbf{A}\) is an integral domain.
\end{enumerate}

For we have \(f(\mathbf{X}) = (\mathbf{X} - \alpha)^p f_1(\mathbf{X}), g(\mathbf{X}) = (\mathbf{X} - \alpha)^q g_1(\mathbf{X})\) with \(f_1(\alpha) \neq 0\) , \(g_1(\alpha) \neq 0\) . Suppose for example that \(p \leqslant q\) ; then we have

\[
f (\mathbf {X}) + g (\mathbf {X}) = (\mathbf {X} - \alpha) ^ {p} \left(f _ {1} (\mathbf {X}) + (\mathbf {X} - \alpha) ^ {q - p} g _ {1} (\mathbf {X})\right),
\]

and if \(p < q\) , \(\alpha\) is not a root of \(f_1(\mathbf{X}) + (\mathbf{X} - \alpha)^{q - p} g_1(\mathbf{X})\) ; this proves (i). On the other hand, we have \(f(\mathbf{X})g(\mathbf{X}) = (\mathbf{X} - \alpha)^{p + q} f_1(\mathbf{X})g_1(\mathbf{X})\) and \(f_1(\alpha)g_1(\alpha) \neq 0\) if \(\mathbf{A}\) is an integral domain; this proves (ii).

PROPOSITION 5. --- Suppose that A is an integral domain. Let f be a non-zero element of A{[}X{]}, and \(\alpha_{1},\ldots,\alpha_{p}\) pairwise distinct roots of f in A, of orders \(k_{1},\ldots,k_{p}\) . We have

\[
f (\mathrm{X}) = \left(\mathrm{X} - \alpha_ {1}\right) ^ {k _ {1}} \left(\mathrm{X} - \alpha_ {2}\right) ^ {k _ {2}} \dots \left(\mathrm{X} - \alpha_ {p}\right) ^ {k _ {p}} g (\mathrm{X})
\]

where \(g \in \mathbf{A}[\mathbf{X}]\) and \(\alpha_1, \ldots, \alpha_p\) are not roots of \(g\) .

We proceed by induction on \(p\) , the proposition being evident for \(p = 1\) , by Def. 1. Suppose then that \(f(\mathbf{X}) = g_1(\mathbf{X})g_2(\mathbf{X})\) , where

\[
g _ {1} (\mathrm{X}) = \left(\mathrm{X} - \alpha_ {1}\right) ^ {k _ {1}} \dots \left(\mathrm{X} - \alpha_ {p - 1}\right) ^ {k _ {p - 1}}, \quad g _ {2} (\mathrm{X}) \in \mathrm{A} [ \mathrm{X} ].
\]

Since A is an integral domain and \(\alpha_{p}\) is distinct from \(\alpha_{1},\ldots,\alpha_{p-1}\) it follows that \(\alpha_{p}\) is not a root of \(g_{1}(X)\) , hence \(\alpha_{p}\) is a root of order \(k_{p}\) of \(g_{2}(X)\) (Prop. 4, (ii)). It follows that \(g_{2}(X)\) is divisible by \((\mathbf{X}-\alpha_{p})^{k_{p}}\) , and so

\[
f (\mathbf {X}) = \left(\mathbf {X} - \alpha_ {1}\right) ^ {k _ {1}} \dots \left(\mathbf {X} - \alpha_ {p}\right) ^ {k _ {p}} g (\mathbf {X})
\]

where \(g(\mathbf{X}) \in \mathbf{A}[\mathbf{X}]\) . Clearly \(\alpha_1, \ldots, \alpha_p\) are not roots of \(g\) .

THEOREM 1. --- Let A be an integral domain. Given a non-zero element f of A{[}X{]}, of degree n, then the sum of the orders of all the roots of f in A is ≤ n. This follows immediately from Prop. 5.

COROLLARY. --- Assume that A is an integral domain and let \(f, g \in A[X]\) , of degrees \(\leqslant n\) . If there exist \(n+1\) pairwise distinct elements \(x_{1}, \ldots, x_{n+1}\) of A such that \(f(x_{i}) = g(x_{i})\) for \(1 \leqslant i \leqslant n+1\) , then \(f = g\) .

It suffices to apply Th. 1 to f - g.

PROPOSITION 6 (Lagrange interpolation formula). --- Let K be a commutative field, \(\alpha_{1}, \alpha_{2}, \ldots, \alpha_{n}\) distinct elements of K and \(\beta_{1}, \beta_{2}, \ldots, \beta_{n}\) any elements of K. For \(i = 1, 2, \ldots, n\) we put

\[
f _ {i} (\mathbf {X}) = \prod_ {j \in U (i)} \left(\mathbf {X} - \alpha_ {j}\right) / \left(\alpha_ {i} - \alpha_ {j}\right),
\]

where \(\mathrm{U}(i)\) is the set of integers j such that \(j \neq i\) and \(1 \leqslant j \leqslant n\) . Then \(\beta_{1}f_{1} + \cdots + \beta_{n}f_{n}\) is the unique element f of K{[}X{]} such that \(\deg f < n\) and \(f(\alpha_{i}) = \beta_{i}\) for \(1 \leqslant i \leqslant n\) .

The uniqueness of \(f\) follows from the Cor. to Th. 1. Let \(f = \beta_{1}f_{1} + \cdots + \beta_{n}f_{n}\) , then since \(f_{i}\) has degree \(n - 1\) , we have \(\deg f < n\) . On the other hand, \(f_{i}(\alpha_{j}) = 0\) for \(j \neq i\) and \(f_{i}(\alpha_{i}) = 1\) , hence \(f(\alpha_{i}) = \beta_{i}\) for \(1 \leqslant i \leqslant n\) .

COROLLARY. --- Suppose that A is an integral domain. Let \(f \in A[X]\) , of degree \textless{} n, and let K be a subring of A which is a field. If there exist n distinct elements \(\alpha_{1}, \ldots, \alpha_{n}\) of A such that \(\alpha_{i} \in K\) and \(f(\alpha_{i}) \in K\) for \(i = 1, \ldots, n\) , then \(f \in K[X]\) .

\subsection{Differential criterion for the multiplicity of a root}

PROPOSITION 7. --- Let \(f \in A[X]\) and let \(\alpha \in A\) be a root of \(f\) . For \(\alpha\) to be a simple root of \(f\) it is necessary and sufficient that \(\alpha\) should not be a root of the derivative \(Df\) of \(f\) .

By hypothesis we have \(f = (\mathbf{X} - \alpha)g\) , where \(g \in \mathbf{A}[\mathbf{X}]\) . For \(\alpha\) to be a simple root of \(f\) it is necessary and sufficient that \(g(\alpha) \neq 0\) . Now we have \(Df = g + (\mathbf{X} - \alpha)Dg\) , whence \((Df)(\alpha) = g(\alpha)\) .

More generally :

PROPOSITION 8. --- Let \(f \in A[X]\) and \(\alpha \in A\) , and suppose that \(\alpha\) has order \(k \geqslant 1\) relative to \(f\) . Then \(\alpha\) has order \(\geqslant k - 1\) relative to \(Df\) . If \(k \cdot 1\) is cancellable in \(A\) , then \(\alpha\) has order \(k - 1\) relative to \(Df\) .

By hypothesis there exists \(g \in \mathbf{A}[\mathbf{X}]\) such that \(f = (\mathbf{X} - \alpha)^k g\) and \(g(\alpha) \neq 0\) . Hence \(Df = k(\mathbf{X} - \alpha)^{k-1}g + (\mathbf{X} - \alpha)^k Dg = (\mathbf{X} - \alpha)^{k-1}(kg + (\mathbf{X} - \alpha)Dg)\) , which establishes the first part of the proposition. The value of \(kg + (\mathbf{X} - \alpha)Dg\) for \(\mathbf{X} = \alpha\) is \(kg(\alpha)\) , and this is non-zero if \(k \cdot 1\) is cancellable in \(\mathbf{A}\) ; this proves the second part of the proposition.

Let \(k\) be an integer \(> 0\) such that \(k \cdot 1 = 0\) in \(\mathbf{A}\) . If \(f(\mathbf{X}) = \mathbf{X}^k\) , then \(0\) is a root of order \(k\) of \(f\) , and a root of arbitrarily high order of \(Df\) .

COROLLARY. --- Let \(f \in \mathbf{A}[\mathbf{X}]\) , \(\alpha \in \mathbf{A}\) and \(p\) an integer \(\geqslant 0\) , further, suppose that \(p!\) . 1 is cancellable in A. Then for \(\alpha\) to be a root of order \(p\) of \(f\) , it is necessary and sufficient that \(\alpha\) should be a root of \(f\) , \(Df\) , \ldots, \(D^{p-1}f\) and not a root of \(D^p f\) .

This follows from Prop. 8 by induction on p.

\subsection{Polynomial functions on an infinite integral domain}

PROPOSITION 9. --- Assume that A is an integral domain. Let I be a set, \((\mathrm{H}_{i})_{i\in I}\) a family of infinite subsets of A and \(H=\prod_{i\in I}H_{i}\subset A^{I}\) . If f is a non-zero

element of \(\mathbf{A}[(\mathbf{X}_i)_{i\in I}]\) , and \(\mathrm{H}_f\) the set of all \(\mathbf{x}\in \mathrm{H}\) such that \(f(\mathbf{x})\neq 0\) , then \(\mathrm{H}\) and \(\mathrm{H}_f\) are equipotent.

\begin{enumerate}
\def\labelenumi{\alph{enumi})}
\tightlist
\item
  First suppose that I is finite and put \(n = \text{Card I}\) . The proposition is clear for \(n = 0\) ; we shall prove it by induction on \(n\) . Choose an element \(i_0\) of I and put \(J = I - \{i_0\}\) , \(B = A[(X_i)_{i \in J}]\) . Since \(f \neq 0\) , we can write \(f = \sum_{k=0}^{m} g_k X_{i_0}^k\) , where \(g_0, \ldots, g_m \in \mathbf{B}\) and \(g_m \neq 0\) . By the induction hypothesis the set \(\mathbf{K}\) of all \(\mathbf{x} \in \prod_{i \in J} \mathrm{H}_i\) such that \(g_m(\mathbf{x}) \neq 0\) is equipotent with \(\prod_{i \in J} \mathrm{H}_i\) . For \(\mathbf{x} \in \mathbf{K}\) the polynomial
\end{enumerate}

\[
h \left(\mathrm{X} _ {i _ {0}}\right) = \sum_ {k = 0} ^ {m} g _ {k} (\mathbf {x}) \mathrm{X} _ {i _ {0}} ^ {k} \in \mathrm{A} \left[ \mathrm{X} _ {i _ {0}} \right]
\]

is non-zero. By Th. 1 (IV, p.~16) the set of \(\alpha \in \mathrm{H}_{i_0}\) such that \(h(\alpha) \neq 0\) is equipotent with \(\mathrm{H}_{i_0}\) whence

\[
\operatorname{Card} \mathrm{H} \geqslant \operatorname{Card} \mathrm{H} _ {f} \geqslant (\operatorname{Card} \mathrm{K}). (\operatorname{Card} \mathrm{H} _ {i _ {0}}) = \operatorname{Card} \mathrm{H},
\]

and so Card \(\mathbf{H} = \mathbf{Card}\mathbf{H}_f\)

\begin{enumerate}
\def\labelenumi{\alph{enumi})}
\setcounter{enumi}{1}
\tightlist
\item
  In the general case there is a finite subset \(\mathbf{I}'\) of \(\mathbf{I}\) such that \(f \in \mathbf{A}[(\mathbf{X}_i)_{i \in \mathbf{I}'}]\) . Let \(\mathbf{H}_f'\) be the set of all \(\mathbf{x} \in \prod_{i \in \mathbf{I}'} \mathbf{H}_i\) such that \(f(\mathbf{x}) \neq 0\) . Then \(\mathbf{H}_f = \mathbf{H}_f' \times \left( \prod_{i \in \mathbf{I} - \mathbf{I}'} \mathbf{H}_i \right)\) , and it suffices to apply the first part of the proof to \(\mathbf{H}_f'\) .
\end{enumerate}

COROLLARY 1. --- We keep the hypothesis and notation of Prop. 9. If I is non-empty, then \(H_{f}\) is infinite.

COROLLARY 2. --- Suppose that A is an infinite integral domain or that A is an algebra over an infinite field. For every \(f \in \mathbf{A}[(\mathbf{X}_i)_{i \in I}]\) , let \(\tilde{f}: \mathbf{A}^{\mathrm{I}} \to \mathbf{A}\) be the polynomial function defined by \(f\) (IV, p.~4). Then the mapping \(f \mapsto \tilde{f}\) is injective.

When \(\mathbf{A}\) is an infinite integral domain, the corollary follows at once from Prop. 9. Suppose that \(\mathbf{A}\) is an algebra over an infinite field \(k\) . Let \(f = \sum_{\nu \in \mathbb{N}^{(I)}} \alpha_{\nu} X^{\nu}\)

be a non-zero element of \(\mathbf{A}[(\mathbf{X}_i)_{i\in I}]\) ; then there exists \(\nu_0\in \mathbf{N}^{(I)}\) such that \(\alpha_{\nu_0}\neq 0\) , and a \(k\) -linear form \(\varphi\) on \(\mathbf{A}\) such that \(\varphi (\alpha_{\nu_0})\neq 0\) . Let \(g = \sum_{\nu \in \mathbf{N}^{(I)}}\varphi (\alpha_{\nu})\mathbf{X}^{\nu}\in k[(\mathbf{X}_i)_{i\in I}]\) ; we have \(g\neq 0\) , hence there exists \(\mathbf{x}\in k^{\mathrm{I}}\) such

\[
g (\mathbf {x}) \neq 0
\]

\[
\varphi (f (\mathbf {x})) = g (\mathbf {x}) \neq 0
\]

\[
f (\mathbf {x}) \neq 0
\]

When A is an infinite integral domain or when A is an algebra over an infinite field, we shall usually identify f with \(\tilde{f}\) .

Suppose that A is finite and let \(f(X) = \prod_{\alpha \in A} (X - \alpha)\) , then \(f \neq 0\) , but \(\tilde{f} = 0\) . For other examples see IV, p.~88, exercises 7 and 8.

THEOREM 2 (Principle of extension of algebraic identities). --- Suppose that A is an infinite integral domain. Let \(g_1, \ldots, g_m\) , \(f\) be elements of \(\mathbf{A}[(\mathbf{X}_i)_{i \in I}]\) and assume the following hypotheses:

\begin{enumerate}
\def\labelenumi{\alph{enumi})}
\item
  \(g_{1} \neq 0, \ldots, g_{m} \neq 0\) ;
\item
  for all \(\mathbf{x} \in \mathbf{A}^{I}\) such that \(g_1(\mathbf{x}) \neq 0, \ldots, g_m(\mathbf{x}) \neq 0\) , we have \(f(\mathbf{x}) = 0\) . Then \(f = 0\) .
\end{enumerate}

For if \(f \neq 0\) , we have \(fg_1 \ldots g_m \neq 0\) (IV, p.~9, Prop. 8), hence there exists \(\mathbf{x} \in \mathbf{A}^{I}\) such that \(f(\mathbf{x})g_1(\mathbf{x}) \ldots g_m(\mathbf{x}) \neq 0\) (IV, p.~18, Cor. 2), which contradicts the hypothesis.

Scholium. --- Let A be an integral domain and \(f \in \mathbf{A}[(\mathbf{X}_i)_{i \in I}]\) . Th. 2 provides a convenient means of proving that f = 0. It suffices to consider an infinite integral domain E containing A as subring; if we can show that \(f((x_i)) = 0\) for all \((x_i) \in \mathrm{E}^{\mathrm{I}}\) (or even for those \((x_i) \in \mathrm{E}^{\mathrm{I}}\) at which a finite number of given non-zero polynomials do not vanish) then it follows that f = 0. If A itself is not infinite, we can for example take E to be the ring A{[}X{]} or its field of fractions.

Once we have proved the relation \(f = 0\) , we can clearly deduce that \(f((y_i)) = 0\) for all \((y_i) \in \mathbf{F}^{I}\) where \(\mathbf{F}\) is any unital associative and commutative \(\mathbf{A}\) -algebra whatsoever; in particular \(\mathbf{F}\) may be finite or non-integral.

In other words, the proof of the identity \(f((x_{i})) = 0\) when the \(x_{i}\) run over an infinite integral domain containing A as subring (with the possible restriction that \(g_{k}((x_{i})) \neq 0\) for \(1 \leqslant k \leqslant m\) , where the \(g_{k}\) are non-zero polynomials) implies the same identity when the \(x_{i}\) run over any unital associative and commutative A-algebra.

In particular let \(f \in \mathbf{Z}[(\mathbf{X}_i)_{i \in I}]\) . If \(f((x_i)) = 0\) when the \(x_i\) run over \(\mathbf{Z}\) (with the possible restriction that \(g_k((x_i)) \neq 0\) for \(1 \leqslant k \leqslant m\) , where the \(g_k\) are non-zero elements of \(\mathbf{Z}[(\mathbf{X}_i)]\) ), then we have the same identity when the \(x_i\) run over an arbitrary commutative ring.

\section{RATIONAL FRACTIONS}

\subsection{Definition of rational fractions}

DEFINITION 1. --- Let K be a commutative field and I a set. The field of fractions (I, p.~116) of the integral domain \(\mathrm{K}[(\mathrm{X}_{i})_{i\in I}]\) is denoted by \(\mathrm{K}((\mathrm{X}_{i})_{i\in I})\) or \(\mathrm{K}(\mathrm{X}_{i})_{i\in I}\) . Its elements are called rational fractions in the indeterminates \(X_{i}\) with coefficients in K.

For \(I = \{1, 2, \dots, n\}\) we write \(K(X_1, X_2, \dots, X_n)\) in place of \(K((X_i)_{i \in I})\) .

Let A be an integral domain and K its field of fractions. The ring \(\mathbf{A}[(\mathbf{X}_{i})_{i\in I}]\) may be identified with a subring of \(\mathbf{K}[(\mathbf{X}_{i})_{i\in I}]\) , hence also of \(\mathbf{K}((\mathbf{X}_{i})_{i\in I})\) . For each \(f\in\mathbf{K}[(\mathbf{X}_{i})_{i\in I}]\) there exists a non-zero element \(\alpha\) of A such that \(\alpha f\in\mathbf{A}[(\mathbf{X}_{i})_{i\in I}]\) . Hence every element of \(\mathbf{K}((\mathbf{X}_{i})_{i\in I})\) can be written as u/v where \(u,v\in\mathbf{A}[(\mathbf{X}_{i})_{i\in I}],v\neq0\) . Thus \(\mathbf{K}((\mathbf{X}_{i})_{i\in I})\) may be identified with the field of fractions of \(\mathbf{A}[(\mathbf{X}_{i})_{i\in I}]\) .

Now let K be a commutative field, I a set and \(J \subset I\) . Put \(\mathbf{B} = \mathbf{K}[(\mathbf{X}_i)_{i \in J}]\) , then \(\mathbf{K}[(\mathbf{X}_i)_{i \in I}] = \mathbf{B}[(\mathbf{X}_i)_{i \in I-J}]\) , and by what has been said, \(\mathbf{K}((\mathbf{X}_i)_{i \in I})\) may be identified with \(\mathbf{K}'((\mathbf{X}_i)_{i \in I-J})\) , where \(\mathbf{K}' = \mathbf{K}((\mathbf{X}_i)_{i \in J})\) .

\subsection{Degrees}

Let K be a commutative field. For every element r of \(\mathrm{K}((\mathbf{X}_{i})_{i\in I})\) there exist u, \(v\in\mathrm{K}[(\mathbf{X}_{i})_{i\in I}]\) such that \(v\neq0\) and \(r=\frac{u}{v}\) . The relation \(\frac{u}{v}=\frac{u_{1}}{v_{1}}\) , where \(v\neq0\) , \(v_{1}\neq0\) is equivalent to \(uv_{1}=vu_{1}\) ; if \(r\neq0\) , we have \(u\neq0\) and \(u_{1}\neq0\) , and so \(\deg u+\deg v_{1}=\deg v+\deg u_{1}\) (IV, p.~9), or also \(\deg u-\deg v=\deg u_{1}-\deg v_{1}\) . The rational integer \(\deg u-\deg v\) thus depends only on r; we call it the degree, or the total degree of r, and denote it by \(\deg r\) . We shall agree to write \(\deg0=-\infty\) . If \(J\subset I\) , we can likewise define the degree with respect to the \(X_{j}\) of index \(j\in J\) . When r is a polynomial, these notions coincide with those defined in IV, p.~2.

PROPOSITION 1. --- Let r, s be two rational fractions.

\begin{enumerate}
\def\labelenumi{(\roman{enumi})}
\tightlist
\item
  If \(\deg r \neq \deg s\) , we have
\end{enumerate}

\[
r + s \neq 0 \quad a n d \quad \deg (r + s) = \sup (\deg r, \deg s).
\]

If \(\deg r = \deg s\) , we have \(\deg (r + s) \leqslant \deg r\) .

\begin{enumerate}
\def\labelenumi{(\roman{enumi})}
\setcounter{enumi}{1}
\tightlist
\item
  We have \(\deg(rs) = \deg r + \deg s\) .
\end{enumerate}

We can limit ourselves to the case where \(r\) and \(s\) are non-zero.

Let us write \(r = \frac{u}{v}\) , \(s = \frac{w}{z}\) , where \(u, v, w, z\) are non-zero polynomials. We have \(rs = \frac{uw}{vz}\) , and so

\[
\deg (r s) = \deg (u w) - \deg (v z) = \deg u - \deg v + \deg w - \deg z = \quad = \deg r + \deg s.
\]

On the other hand, we have \(r + s = \frac{uz + vw}{vz}\) . Suppose that \(\deg r \neq \deg s\) , in other words \(\deg u + \deg z \neq \deg w + \deg v\) . Then \(uz + wv \neq 0\) , and

\[
\begin{array}{r l} \deg (r + s) & = \deg (u z + w v) - \deg (v z) \\ & = \sup (\deg (u z), \deg (w v)) - \deg (v z) \\ & = \sup (\deg (u z) - \deg (v z), \deg (w v) - \deg (v z)) \\ & = \sup (\deg r, \deg s). \end{array}
\]

Suppose that \(\deg r = \deg s\) , that is \(\deg u + \deg z = \deg w + \deg v\) . If \(r + s \neq 0\) , then we have

\[
\begin{array}{r l} \deg (r + s) & = \deg (u z + w v) - \deg (v z) \\ & \leqslant \deg (u z) - \deg (v z) = \deg r. \end{array}
\]

* The mapping \(r \mapsto -\deg r\) is thus a discrete valuation on the field \(K((X_i)_{i \in I})\) . *

\subsection{Substitutions}

Let K be a commutative field, E a unital associative K-algebra, \(\mathbf{x} = (x_i)_{i \in I}\) a family of pairwise permutable elements of E. Let \(\mathbf{B} = \mathbf{K}[(\mathbf{X}_i)_{i \in I}]\) and \(S_x\) the set of all non-zero \(v \in B\) such that \(v(\mathbf{x})\) is invertible in E. Let \(u \in B\) , \(v \in S_x\) and \(f = \frac{u}{v} \in \mathbf{K}((\mathbf{X}_i)_{i \in I})\) . The element \(u(\mathbf{x}) v(\mathbf{x})^{-1} = v(\mathbf{x})^{-1} u(\mathbf{x})\) is defined in E; moreover, if \(u_1\) , \(v_1\) are two polynomials such that \(f = \frac{u_1}{v_1}\) and \(v_1 \in S_x\) , then \(uv_1 = vu_1\) , hence \(u(\mathbf{x}) v_1(\mathbf{x}) = v(\mathbf{x}) u_1(\mathbf{x})\) and so

\[
u (\mathbf {x}) v (\mathbf {x}) ^ {- 1} = u _ {1} (\mathbf {x}) v _ {1} (\mathbf {x}) ^ {- 1}.
\]

Let \(f \in \mathbf{K}((\mathbf{X}_i)_{i \in I})\) . If there exists at least one couple \((u, v)\) such that \(f = \frac{u}{v}\) and \(v \in S_x\) , we shall say that \(x\) is substitutable in \(f\) ; the element \(u(x) v(x)^{-1}\) which only depends on \(f\) and \(x\) is then denoted by \(f(x)\) or \(f((x_i))\) or \(f((x_i)_{i \in I})\) .

PROPOSITION 2. --- Let K be a commutative field, E a unital associative K-algebra and \(\mathbf{x} = (x_i)_{i \in I}\) a family of pairwise permutable elements of E. The set \(\mathbf{S}_{\mathbf{x}}^{-1}\mathbf{B}\) of \(f \in \mathbf{K}((\mathbf{X}_i)_{i \in I})\) such that \(\mathbf{x}\) is substitutable in \(f\) is a K-subalgebra of \(\mathbf{K}((\mathbf{X}_i)_{i \in I})\) . The mapping \(f \mapsto f(\mathbf{x})\) is a unital homomorphism \(\varphi\) of \(\mathbf{S}_{\mathbf{x}}^{-1}\mathbf{B}\) into E. The image \(\varphi(\mathbf{S}_{\mathbf{x}}^{-1}\mathbf{B})\) is the set of all \(yz^{-1}\) , where \(y\) runs over the unital subalgebra \(\mathbf{K}[\mathbf{x}]_E\) of E generated by the family \(\mathbf{x}\) and where \(z\) runs over the set of all invertible elements of \(\mathbf{K}[\mathbf{x}]_E\) .

Let \(f_{1} = \frac{u_{1}}{v_{1}}\) , \(f_{2} = \frac{u_{2}}{v_{2}}\) be two elements of \(K((X_{i})_{i\in I})\) such that \(v_{1}, v_{2} \in S_{x}\) . We have \(f_{1} + f_{2} = \frac{u_{1}v_{2} + u_{2}v_{1}}{v_{1}v_{2}}\) , \(f_{1}f_{2} = \frac{u_{1}u_{2}}{v_{1}v_{2}}\) and \(v_{1}, v_{2} \in S_{x}\) . Hence \(S_{x}^{-1}B\) is a K-subalgebra of \(K((X_{i})_{i\in I})\) . The rest of the proposition is clear.

COROLLARY. --- Let L be a commutative field, K a subfield of L, \(\mathbf{x} = (x_i)_{i \in I}\) a family of elements of L, M the set consisting of the \(x_i\) , U the set of all \(f \in \mathrm{K}((\mathrm{X}_i)_{i \in I})\) such that x is substitutable in f and \(\varphi\) the homomorphism \(f \mapsto f(\mathbf{x})\) of U into L. Then \(\varphi(\mathrm{U})\) is the subfield of L generated by \(K \cup M\) .

Let \(L'\) be the subfield of L generated by \(K \cup M\) . We have

\[
\mathbf {K} \cup \mathbf {M} \subset \varphi (\mathbf {U}) \subset \mathbf {L} ^ {\prime}
\]

and \(\varphi(\mathbf{U})\) is a subring of \(\mathbf{L}\) . Now Prop. 2 implies that \(\varphi(\mathbf{U})\) is a subfield of \(\mathbf{L}\) , whence \(\varphi(\mathbf{U}) = \mathbf{L}'\) .

Let \(f \in \mathbf{K}((\mathbf{X}_i)_{i \in I})\) and let \((g_i)_{i \in I}\) be a family of elements of \(\mathbf{K}((\mathbf{Y}_l)_{l \in L})\) . If \((g_{i})\) is substitutable in f, then \(f((g_{i}))\) is an element of \(\mathbf{K}((\mathbf{Y}_{l})_{l\in L})\) . In particular, \((\mathbf{X}_{i})_{i\in I}\) is substitutable in f and \(f=f((\mathbf{X}_{i})_{i\in I})\) .

PROPOSITION 3. --- Let E be an algebra over K which is associative, commutative, unital and non-zero. Let \(f \in \mathbf{K}((\mathbf{X}_i)_{i \in \mathbb{I}})\) and for each \(i \in \mathbb{I}\) , let \(g_i \in \mathbf{K}((\mathbf{Y}_l)_{l \in \mathbb{I}})\) . Given a family \(\mathbf{y} = (y_l)_{l \in \mathbb{L}}\) of elements of E, suppose that \(\mathbf{y}\) is substitutable into each \(g_i\) and \((g_i(\mathbf{y}))_{i \in \mathbb{I}}\) is substitutable in \(f\) . Then:

\begin{enumerate}
\def\labelenumi{(\roman{enumi})}
\item
  \((g_{i})_{i\in I}\) is substitutable in \(f\) ;
\item
  if we denote by \(h\) the element \(f((g_i))\) of \(\mathbf{K}((\mathbf{Y}_l)_{l \in L})\) , then \(\mathbf{y}\) is substitutable in \(h\) and \(h(\mathbf{y}) = f((g_i(\mathbf{y})))\) .
\end{enumerate}

We may take I to be finite. By hypothesis, for each \(i \in I\) , \(g_i\) can be put in the form \(p_i/q_i\) where \(p_i\) , \(q_i \in K[(Y_l)_{l \in L}]\) and \(q_i(\mathbf{y})\) is invertible in E. Likewise \(f\) can be written as \(u/v\) , where \(u\) , \(v \in K[(X_i)_{i \in I}]\) and \(v((g_i(\mathbf{y})))\) is invertible. Let \(m = \sup(\deg u, \deg v)\) , and let \(w = \prod_{i \in I} q_i \in K[(Y_l)_{l \in L}]\) , \(u_1 = u((g_i))w^m\) , \(v_1 = v((g_i))w^m\) . The polynomial \(u\) is a K-linear combination of monomials \(\prod_{i \in I} X_i^{\nu_i}\) such that \(\sum_{i \in I} \nu_i \leqslant m\) . We have \(w^m \prod_{i \in I} g_i^{\nu_i} = w^m\left(\prod_{i \in I} p_i^{\nu_i}\right)\left(\prod_{i \in I} q_i^{\nu_i}\right)^{-1} \in K[(Y_l)_{l \in L}]\) by the choice of \(m\) . Hence \(u_1 \in K[(Y_l)_{l \in L}]\) and similarly \(v_1 \in K[(Y_l)_{l \in L}]\) . Moreover, \(v_1(\mathbf{y}) = (w(\mathbf{y}))^m v((g_i(\mathbf{y})))\) is invertible. Hence \(v_1 \neq 0\) , because \(E \neq 0\) , and so \(v((g_i)) \neq 0\) . The family \((g_i)\) is thus substitutable in \(f\) . Besides we have \(f((g_i)) = u_1/v_1\) , hence \(\mathbf{y}\) is substitutable in \(h = f((g_i))\) , and \(h(\mathbf{y}) = u_1(\mathbf{y})/v_1(\mathbf{y}) = u((g_i(\mathbf{y})))/v((g_i(\mathbf{y}))) = f((g_i(\mathbf{y})))\) .

Let K be a commutative field, E a commutative associative and unital K-algebra, and let \(f \in \mathbf{K}((\mathbf{X}_i)_{i \in I})\) . Let \(T_f\) be the set of all \(\mathbf{x} = (x_i)_{i \in I} \in E^{I}\) which are substitutable in \(f\) . The mapping \(\mathbf{x} \mapsto f(\mathbf{x})\) of \(T_f\) into E is called the rational function associated with \(f\) (and E); we sometimes denote it by \(\tilde{f}\) . If \(g \in \mathbf{K}((\mathbf{X}_i)_{i \in I})\) we have \(T_f \cap T_g \subset T_{f + g}\) , \(T_f \cap T_g \subset T_{fg}\) , hence the rational function associated with \(f + g\) (resp. \(fg\) ) is defined on \(T_f \cap T_g\) and has the same value on this set as \(\tilde{f} + \tilde{g}\) (resp. \(\tilde{f}\tilde{g}\) ). Let \(T_f'\) be the set of \(\mathbf{x} \in T_f\) such that \(f(\mathbf{x})\) is invertible; if \(\mathbf{x} \in T_f'\) , \(\mathbf{x}\) is substitutable in \(1/f\) and the rational function associated with \(1/f\) takes at \(\mathbf{x}\) the value \(f(\mathbf{x})^{-1}\) .

If K is an infinite commutative field, \(f \in \mathbf{K}((\mathbf{X}_i)_{i \in I})\) , \(g \in \mathbf{K}((\mathbf{X}_i)_{i \in I})\) and \(\tilde{f}\) , \(\tilde{g}\) are the rational functions associated with f, g (and K), and if \(\tilde{f}(\mathbf{x}) = \tilde{g}(\mathbf{x})\) for all \(x \in T_f \cap T_g\) then f = g. For if f = u/v and \(g = u_1 / v_1\) , where u, v, \(u_1\) , \(v_1\) are polynomials, we have \(u(\mathbf{x})v_1(\mathbf{x}) = u_1(\mathbf{x})v(\mathbf{x})\) for all x such that \(v(x)v_1(x) \neq 0\) , hence \(uv_1 = u_1v\) (IV, p.~18, Th. 2). Therefore the mapping \(f \mapsto \tilde{f}\) is injective and we shall often identify f and \(\tilde{f}\) .

* Using the factoriality of \(\mathbf{K}[(\mathbf{X}_i)_{i\in I}]\) (Comm. Alg., VII, § 3, No.~2 p.~502 and Cor. of Th. 2 p.~506), one easily shows the following: for every \(f\in \mathbf{K}((\mathbf{X}_i)_{i\in I})\) there exist \(u,v\in \mathbf{K}[(\mathbf{X}_i)_{i\in I}]\) such that:

\begin{enumerate}
\def\labelenumi{\arabic{enumi})}
\item
  \(f = u / v\)
\item
  for \(\mathbf{x} \in \mathbf{K}^{I}\) to be substitutable in \(f\) it is necessary and sufficient that \(v(\mathbf{x}) \neq 0\) .
\end{enumerate}

\subsection{Differentials and derivations}

Let K be a commutative field. By III, p.~558, Prop. 5, every derivation D of \(\mathbf{K}[(\mathbf{X}_i)_{i\in I}]\) extends in a unique fashion to a derivation \(\bar{\mathbf{D}}\) of \(\mathbf{K}((\mathbf{X}_i)_{i\in I})\) . If D, D' are permutable derivations of \(\mathbf{K}[(\mathbf{X}_i)_{i\in I}]\) , then the bracket \([\mathbf{D},\mathbf{D}'] = \mathbf{DD}' - \mathbf{D}'\mathbf{D}\) is zero, hence \([\bar{\mathbf{D}},\bar{\mathbf{D}}']\) which is a derivation of \(\mathbf{K}((\mathbf{X}_i)_{i\in I})\) extending {[}D, D'{]} is zero; in other words, \(\bar{\mathbf{D}}\) and \(\bar{\mathbf{D}}'\) are permutable. In particular the derivations D\_i (IV, p.~6) extend to derivations of \(\mathbf{K}((\mathbf{X}_i)_{i\in I})\) again denoted by D\_i and which are pairwise permutable. If \(f\in \mathbf{K}((\mathbf{X}_i)_{i\in I})\) , D\_if is also written D\_x\_i f or \(\frac{\partial f}{\partial\mathbf{X}_i}\) or \(f_{\mathbf{x}_i}'\) . When there is only a single indeterminate X one uses the notation Df, \(\frac{df}{d\mathbf{X}}\) , \(f'\) .

Let \(\mathbf{B} = \mathbf{K}[(\mathbf{X}_i)_{i \in I}]\) , \(\mathbf{C} = \mathbf{K}((\mathbf{X}_i)_{i \in I})\) . By III, p.~574, Prop. 23, the canonical mapping

\[
\Omega_ {\mathrm{K}} (\mathrm{B}) \otimes_ {\mathrm{B}} \mathrm{C} \rightarrow \Omega_ {\mathrm{K}} (\mathrm{C})
\]

is an isomorphism of vector C-spaces. Bearing in mind III, p.~570, we see that the vector C-space \(\Omega_{\mathbf{K}}(\mathbf{C})\) admits as a basis the family \((d\mathbf{X}_i)_{i\in I}\) of the differentials of the \(\mathbf{X}_i\) . Let \(\partial_i\) be the coordinate form of index \(i\) on \(\Omega_{\mathbf{K}}(\mathbf{C})\) relative to that basis. Then the mapping \(u\mapsto \langle \partial_i,du\rangle\) of C into itself is a derivation of C which maps \(\mathbf{X}_i\) to 1 and \(\mathbf{X}_j\) to 0 for \(j\neq i\) , and so is equal to \(\mathrm{D}_i\) ; in other words, we have

\[
d u = \sum_ {i \in I} \left(\mathrm{D} _ {i} u\right) d \mathrm{X} _ {i}\tag{1}
\]

for every \(u \in \mathbf{C}\) . If \(\mathbf{I}\) is finite, \((\mathrm{D}_i)_{i \in \mathbf{I}}\) is a basis of the vector \(\mathbf{C}\) -space of derivations of \(\mathbf{C}\) .

PROPOSITION 4. --- Let E be an associative, commutative and unital K-algebra, \(\mathbf{x} = (x_i)_{i \in I}\) a family of elements of E and \(f \in \mathbf{K}((\mathbf{X}_i)_{i \in I})\) . Suppose that \(\mathbf{x}\) is substitutable in \(f\) and \(y = f(\mathbf{x})\) .

\begin{enumerate}
\def\labelenumi{(\roman{enumi})}
\item
  For every derivation \(\Delta\) of \(\mathbf{K}((\mathbf{X}_i)_{i\in I})\) which maps \(\mathbf{K}[(\mathbf{X}_i)_{i\in I}]\) into itself, \(\mathbf{x}\) is substitutable in \(\Delta f\) .
\item
  For every derivation D of E into an E-module we have
\end{enumerate}

\[
\mathrm{D} y = \sum_ {i \in I} (\mathrm{D} _ {i} f) (\mathbf {x}) \cdot \mathrm{D} x _ {i}.
\]

Let \(f = \frac{u}{v}\) with \(u, v \in \mathbf{K}[(\mathbf{X}_i)_{i \in I}]\) and \(v(\mathbf{x})\) invertible in E. Let \(\Delta\) be a derivation of \(\mathbf{K}((\mathbf{X}_i)_{i \in I})\) mapping \(\mathbf{K}[(\mathbf{X}_i)_{i \in I}]\) into itself. We have

\[
\Delta f = \frac {(\Delta u) v - u (\Delta v)}{v ^ {2}}
\]

and \(v^2(\mathbf{x})\) is invertible, hence \(\mathbf{x}\) is substitutable in \(\Delta f\) . Secondly put \(r = u(\mathbf{x})\) , \(s = v(\mathbf{x})\) ; we have \(y = s^{-1}r\) , hence for every derivation D of E into an E-module we have

\[
\begin{array}{l} \mathrm{D} y = s ^ {- 2} (s (\mathrm{D} r) - r (\mathrm{D} s)) \\ = s ^ {- 2} \left(s \sum_ {i \in I} (\mathrm{D} _ {i} u) (\mathbf {x}) \cdot \mathrm{D} x _ {i} - r \sum_ {i \in I} (\mathrm{D} _ {i} v) (\mathbf {x}) \cdot \mathrm{D} x _ {i}\right) \end{array}
\]

by Prop. 4 of IV, p.~6. Thus \(Dy = \sum_{i \in I} w_i \cdot Dx_i\) with

\[
w _ {i} = v (\mathbf {x}) ^ {- 2} (v (\mathbf {x}) \left(\mathrm{D} _ {i} u\right) (\mathbf {x}) - u (\mathbf {x}) \left(\mathrm{D} _ {i} v\right) (\mathbf {x})) = \left(\mathrm{D} _ {i} f\right) (\mathbf {x}).
\]

\section{FORMAL POWER SERIES}

\subsection{Definition of formal power series. Order}

Let I be a set. We recall (III, p.~454 and 456) that the total algebra of the monoid \(\mathbf{N}^{(I)}\) over A is called the algebra of formal power series with respect to the indeterminates \(\mathbf{X}_i\) ( \(i \in I\) ) (or in the indeterminates \(\mathbf{X}_i\) ) with coefficients in A. It is denoted by \(\mathbf{A}[[\mathbf{X}_i]]_{i \in I}\) or \(\mathbf{A}[[(\mathbf{X}_i)_{i \in I}]]\) or also \(\mathbf{A}[[\mathbf{X}]]\) , on denoting by \(\mathbf{X}\) the family \((\mathbf{X}_i)_{i \in I}\) : in this paragraph we shall mainly use the notation \(\mathbf{A}[[I]]\) . Sometimes it is convenient to designate the canonical image in \(\mathbf{A}[[I]]\) of the element \(i\) of I by a symbol other than \(\mathbf{X}_i\) , for example \(\mathbf{Y}_i\) , \(\mathbf{Z}_i\) , \(\mathrm{T}_i\) , \ldots; the conventions used in this case are analogous to those for polynomials (IV, p.~1). The algebra \(\mathbf{A}[[I]]\) is then designated by \(\mathbf{A}[[\mathbf{Y}_i]]_{i \in I}\) , or \(\mathbf{A}[[\mathbf{Y}]]\) etc.

When I is a finite set of p elements, we also say that A{[}{[}I{]}{]} is an algebra of formal power series in p indeterminates. These algebras are all isomorphic, for fixed p.~An algebra of formal power series in 1, 2, \ldots{} indeterminates will also be denoted by A{[}{[}X{]}{]}, A{[}{[}U, V{]}{]}, \ldots, the set I of indices not being specified.

A formal power series \(u\) is conventionally written \(u = \sum_{\nu \in \mathbb{N}^{(I)}} \alpha_{\nu} X^{\nu}\) (cf.~IV, p.~1).

The \(\alpha_{\nu}\) are the coefficients of \(u\) ; there may be infinitely many of them \(\neq 0\) . The \(\alpha_{\nu}X^{\nu}\) are called terms of \(u\) ; for \(u\) to be a polynomial it is necessary and sufficient that \(u\) should have only a finite number of terms \(\neq 0\) . The terms \(\alpha_{\nu}X^{\nu}\) such that \(|\nu|=p\) are called the terms of total degree p.~The formal power series \(u_{p}=\sum_{|\nu|=p}\alpha_{\nu}X^{\nu}\) is called the homogeneous component of degree p of u (it is a polynomial when I is finite); \(u_{0}\) is identified with an element of A called also the constant term of u. We say that u is homogeneous of degree p if \(u = u_{p}\) . If u, \(v \in A[[I]]\) and w = uv, we have

\[
w _ {p} = \sum_ {q + r = p} u _ {q} v _ {r}\tag{1}
\]

for every integer \(p \geqslant 0\) .

We recall (III, p.~456), that the order \(\omega(u)\) of a formal power series \(u \neq 0\) is the least integer \(p\) such that \(u_p \neq 0\) . We shall agree to adjoin to \(\mathbf{Z}\) an element written \(\infty\) and extend the order relation and addition from \(\mathbf{Z}\) to \(\mathbf{Z} \cup \{\infty\}\) by the conventions

\[
n <   \infty , \quad \infty + \infty = \infty , \quad \infty + n = n + \infty = \infty
\]

for every \(n \in \mathbf{Z}\) ; we also put \(\omega(0) = \infty\) . With these conventions we have the relations

\[
\begin{array}{c} \omega (u + v) \geqslant \inf (\omega (u), \omega (v)), \\ \omega (u + v) = \inf (\omega (u), \omega (v)) \quad \text { if } \quad \omega (u) \neq \omega (v), \\ \omega (u v) \geqslant \omega (u) + \omega (v), \end{array}
\]

for any formal power series \(u\) and \(v\) in A{[}{[}I{]}{]}.

We recall (III, p.~457) that for any subset \(J\) of \(I\) we identify \(A[[I]]\) with \(A[[I - J]][[J]]\) , which allows us to define the order \(\omega_{j}(u)\) of a formal power series with respect to the \(X_{j}\) ( \(j \in J\) ), the homogeneous component of \(u\) with respect to the \(X_{j}\) ( \(j \in J\) ) etc.

Let \(\varphi\) be a homomorphism of A into a ring B. We extend \(\varphi\) to a homomorphism \(\overline{\varphi}\) of A{[}{[}I{]}{]} into B{[}{[}I{]}{]} by letting each formal power series \(u = \sum_{\nu} \alpha_{\nu} X^{\nu}\) correspond

to the formal power series \(\sum_{\nu} \varphi(\alpha_{\nu}) X^{\nu}\) ; we say that the latter is obtained by

applying \(\varphi\) to the coefficients of the formal power series \(u\) . We shall sometimes write \({}^{\varphi}u\) for \(\bar{\varphi}(u)\) .

In particular if A is a subring of B and \(\varphi\) the canonical injection of A into B, then the homomorphism \(\bar{\varphi}\) of A{[}{[}I{]}{]} into B{[}{[}I{]}{]} is injective; we shall in general identify A{[}{[}I{]}{]} with a subring of B{[}{[}I{]}{]} by means of \(\bar{\varphi}\) .

\subsection{Topology on the set of formal power series. Summable families}

By definition A{[}{[}I{]}{]} is nothing other than the product set \(\mathbf{A}^{\mathbf{N}^{(I)}}\) . Except for express mention to the contrary we shall equip A with the discrete topology and A{[}{[}I{]}{]} with the product topology (Gen.~Top. I, p.~31 f.) which we shall call the canonical topology. Equipped with addition and the discrete topology, A is a separated and complete topological group; hence for addition A{[}{[}I{]}{]} is a separated and complete topological group (Gen.~Top., III, p.~238 and 242 and Gen.~Top., II, p.~187). Moreover the algebra A{[}(X \(_{i}\) ) \(_{i\in I}\) {]} of polynomials is dense in A{[}{[}I{]}{]} (Gen.~Top., III p.~238, Prop. 25) and we may thus consider A{[}{[}I{]}{]} as the completion of A{[}(X \(_{i}\) ) \(_{i\in I}\) {]}.

For each \(\beta \in \mathbf{N}^{(I)}\) let \(S_{\beta}\) be the set of multi-indices \(\nu\) such that \(\nu \leqslant \beta\) and let \(\alpha_{\beta}\) be the set of formal power series \(u = \sum_{\nu} \alpha_{\nu} X^{\nu}\) such that \(\alpha_{\nu} = 0\) for

\(\nu \in S_{\beta}\) . Clearly \(S_{\beta}\) is a finite subset of \(\mathbf{N}^{(I)}\) , and every finite subset of \(\mathbf{N}^{(I)}\) is contained in a set of the form \(S_{\beta}\) . It follows that the family \((\mathfrak{a}_{\beta})_{\beta \in \mathbf{N}^{(I)}}\) is a fundamental system of neighbourhoods of 0 in A{[}{[}I{]}{]}. The sets \(\mathfrak{a}_{\beta}\) are ideals in A{[}{[}I{]}{]}, hence (Gen.~Top., III, p.~275) A{[}{[}I{]}{]} is a topological ring.

Lemma 1. --- Let L be an infinite set and \((u_{\lambda})_{\lambda \in L}\) a family of elements of A{[}{[}I{]}{]}, and put \(u_{\lambda} = \sum_{\nu} \alpha_{\lambda, \nu} X^{\nu}\) for \(\lambda \in L\) . Then the following conditions are

equivalent :

\begin{enumerate}
\def\labelenumi{(\roman{enumi})}
\item
  The family \((u_{\lambda})_{\lambda \in L}\) is summable (Gen.~Top., III, p.~262) in A{[}{[}I{]}{]}.
\item
  We have \(\lim u_{\lambda} = 0\) , taken along the filter of complements of finite subsets of \(L\) .
\item
  For every \(\nu \in \mathbf{N}^{(I)}\) we have \(\alpha_{\lambda, \nu} = 0\) except for a finite number of indices \(\lambda \in L\) .
\end{enumerate}

When these conditions hold, the series \(u = \sum_{\lambda \in L} u_{\lambda}\) is equal to \(\sum_{\nu} \alpha_{\nu} X^{\nu}\) with \(\alpha_{\nu} = \sum_{\lambda \in L} \alpha_{\lambda, \nu}\) for each \(\nu \in \mathbf{N}^{(I)}\) .

The equivalence of (i) and (ii) follows from Cor. 2 of Gen.~Top., III, p.~263. The equivalence of (ii) and (iii) follows from the properties of limits in a product space (Gen.~Top., I, p.~55, Cor. 1).

The last assertion follows from Prop. 4 of Gen.~Top., III, p.~266.

Let us give some examples of summable families.

\begin{enumerate}
\def\labelenumi{\alph{enumi})}
\item
  Let \(u \in \mathbf{A}[[\mathrm{I}]]\) and let \(\alpha_{\nu}\) be the coefficient of \(\mathbf{X}^{\nu}\) in \(u\) . The family \((\alpha_{\nu}\mathbf{X}^{\nu})_{\nu \in \mathbb{N}^{(I)}}\) is then summable, with sum \(u\) (which justifies writing \(u = \sum_{\nu} \alpha_{\nu}\mathbf{X}^{\nu}\) ).
\item
  Let \(u \in \mathbf{A}[[\mathbf{I}]]\) ; for every integer \(p \geqslant 0\) let \(u_p\) be the homogeneous component of degree \(p\) of \(u\) . Then the family \((u_p)_{p \in \mathbb{N}}\) is summable and we have \(u = \sum_{p \geqslant 0} u_p\) .
\item
  Let \((u_{\lambda})_{\lambda \in L}\) be a family of elements of A{[}{[}I{]}{]} and suppose that for every integer \(n \geqslant 0\) the set of \(\lambda \in L\) such that \(\omega(u_{\lambda}) < n\) is finite. Then the family \((u_{\lambda})_{\lambda \in L}\) is summable.
\end{enumerate}

Remark. --- Suppose that I is finite. For every integer \(n \geqslant 0\) let \(\mathbf{b}_n\) be the set of formal power series \(u \in \mathbf{A}[[\mathbf{I}]]\) such that \(\omega(u) \geqslant n\) . The sequence \((\mathbf{b}_n)_{n \geqslant 0}\) is a fundamental system of neighbourhoods of 0 in A{[}{[}I{]}{]}. Therefore a family of elements \(u_{\lambda}\) of A{[}{[}I{]}{]} ( \(\lambda \in L\) ) is summable if and only if for every \(n \in \mathbb{N}\) the set of \(\lambda \in L\) such that \(\omega(u_{\lambda}) < n\) is finite.

PROPOSITION 1. --- Let \((u_{\lambda})_{\lambda \in L}\) and \((v_{\mu})_{\mu \in M}\) be two summable families of elements of A{[}{[}I{]}{]}. Then the family \((u_{\lambda}v_{\mu})_{(\lambda, \mu) \in L \times M}\) is summable and we have

\[
\sum_ {(\lambda , \mu) \in L \times M} u _ {\lambda} v _ {\mu} = \left(\sum_ {\lambda \in L} u _ {\lambda}\right) \left(\sum_ {\mu \in M} v _ {\mu}\right).\tag{2}
\]

Let \((\alpha_{\lambda, \nu})_{\nu \in \mathbf{N}^{(I)}}\) (resp. \((\beta_{\mu, \nu})_{\nu \in \mathbf{N}^{(I)}}\) ) be the family of coefficients of \(u_{\lambda}\) (resp. \(v_{\mu}\) ). For each \(\nu \in \mathbf{N}^{(I)}\) there exists only a finite number of pairs \((\nu_1, \nu_2) \in \mathbf{N}^{(I)} \times \mathbf{N}^{(I)}\) such that \(\nu_1 + \nu_2 = \nu\) , hence only a finite number of pairs \((\lambda, \mu) \in L \times M\) such that the coefficient of \(X^{\nu}\) in \(u_{\lambda}v_{\mu}\) is \(\neq 0\) . Hence the family \((u_{\lambda}v_{\mu})_{(\lambda, \mu) \in L \times M}\) is summable. Now the formula (2) follows from the associativity of the sum (Gen.~Top., III, p.~265, formula (2)).

In A{[}{[}I{]}{]} the product is an associative and commutative composition law. We may therefore speak of a multipliable family of elements of A{[}{[}I{]}{]} and of the product of a multipliable family (Gen.~Top., III, p.~262, remark 3).

PROPOSITION 2. --- Let \((u_{\lambda})_{\lambda \in L}\) be a summable family of elements of A{[}{[}I{]}{]}. (i) The family \((1 + u_{\lambda})_{\lambda \in L}\) is multipliable.

\begin{enumerate}
\def\labelenumi{(\roman{enumi})}
\setcounter{enumi}{1}
\tightlist
\item
  Let \(\mathfrak{F}\) be the set of all finite subsets of L. For any \(\mathbf{M} \in \mathfrak{F}\) put \(u_{\mathbf{M}} = \prod_{\lambda \in \mathbf{M}} u_{\lambda}\) . Then the family \((u_{\mathbf{M}})_{\mathbf{M} \in \mathfrak{F}}\) is summable and we have
\end{enumerate}

\[
\sum_ {M \in \mathfrak {F}} u _ {M} = \prod_ {\lambda \in L} (1 + u _ {\lambda}).
\]

Let us define the ideals \(\mathfrak{a}_{\beta}\) as at the beginning of this No., and let \(\beta \in \mathbf{N}^{(I)}\) . There exists a finite subset \(L_0\) of \(L\) such that \(u_{\lambda} \in \mathfrak{a}_{\beta}\) for \(\lambda \notin L_0\) . Then for every \(M \in \mathfrak{F}\) such that \(M \subset L_0\) we have \(u_M \in \mathfrak{a}_{\beta}\) . It follows that the family \((u_M)_{M \in \mathfrak{F}}\) is summable. On the other hand, for any finite subset \(M_0\) of \(L\) we have

\[
\sum_ {M \subset M _ {0}} u _ {M} = \prod_ {\lambda \in M _ {0}} (1 + u _ {\lambda}).
\]

Taken along the filtered ordered set \(\mathfrak{F}\) , the left-hand side has as limit \(\sum_{M\in\mathfrak{F}}u_{M}\) . Hence the right-hand side has as limit \(\sum_{M\in\mathfrak{F}}u_{M}\) , which proves (i) and (ii) at

the same time.

PROPOSITION 3. --- Let \(u = \sum_{\nu} \alpha_{\nu} X^{\nu} \in \mathbf{A}[[\mathbf{I}]]\) and \(m\) an integer \(> 0\) . For every \(n \in \mathbb{N}\) let \((\alpha_{\nu, n})_{\nu \in \mathbb{N}^{(I)}}\) be the family of coefficients of \(u^n\) . If \(\alpha_0^m = 0\) , then \(\alpha_{\nu, n} = 0\) for \(n \geqslant |\nu| + m\) .

Let \(\nu \in \mathbf{N}^{(I)}\) and \(n\in \mathbf{N}\) . We have

\[
\alpha_ {\nu , n} = \sum_ {\nu (1) + \dots + \nu (n) = \nu} \alpha_ {\nu (1)} \dots \alpha_ {\nu (n)}.
\]

If \(n \geqslant |\nu| + m\) and \(\nu(1) + \cdots + \nu(n) = \nu\) , we have \(|\nu(1)| + \cdots + |\nu(n)| \leqslant n - m\) . We thus have \(\nu(r) = 0\) and so \(\alpha_{\nu(r)} = \alpha_0\) for at least m distinct values of r; it follows that \(\alpha_{\nu(1)} \ldots \alpha_{\nu(n)} = 0\) , whence the result.

COROLLARY. --- Let \(u \in \mathbf{A}[[\mathbf{I}]]\) ; then for \(\lim_{n \to \infty} u^n = 0\) to hold it is necessary and

sufficient that the constant term of \(u\) should be nilpotent.

Let \(\alpha_0\) be the constant term of \(u\) . The constant term of \(u^n\) is \(\alpha_0^n\) , hence the stated condition is necessary; it is sufficient by Prop. 3.
